% ============================================================
% Corpo das questões — GERADO por gerar-corpo.py a partir de
% conteudo-blocos.json (SSOT). Não editar à mão: editar o JSON
% e regenerar. Sincronia prova/solução garantida via \ifsolucao
% ============================================================

\ifsolucao\else
\begin{center}
\renewcommand{\arraystretch}{1.3}
\begin{tabular}{|l|l|c|c|}
\hline
\textbf{Bloco} & \textbf{Conteúdo} & \textbf{Aulas} & \textbf{Pontos} \\ \hline
1 & Teoria do Consumidor & 1--3 & 25 \\ \hline
2 & Equilíbrio Geral & 4--6 & 25 \\ \hline
3 & Externalidades e Bens Públicos & 7 & 25 \\ \hline
4 & Assimetria Informacional e Sinalização & 8--9 & 25 \\ \hline
\multicolumn{3}{|r|}{\textbf{Total}} & \textbf{100} \\ \hline
\end{tabular}
\end{center}
\vspace{0.5em}
\fi

\blocotitulo{Bloco 1 --- Teoria do Consumidor \hfill \normalsize (25 pontos)}

\textbf{Questão 1.} (10 pontos) \textbf{Itens objetivos.} Nos itens \textbf{a)} e \textbf{b)}, assinale a \emph{única} alternativa correta (3 pontos cada; não há penalização por erro). Nos itens \textbf{c)} e \textbf{d)}, responda \textbf{Verdadeiro ou Falso} (2 pontos cada; vale a regra de anulação descrita na capa). Registre suas respostas no quadro-resposta ao final do bloco.

\textbf{a)} (3 pontos) Uma operadora (suponha a Vivo) oferece pacotes em que cada GB de dados ($x_1$) só é útil acompanhado de um bloco de minutos de voz ($x_2$), na proporção fixa de $1$ para $1$, de modo que as preferências do assinante são $u(x_1,x_2)=\min\{x_1,x_2\}$. Suponha que o assinante disponha de $m=R\$~60$ e enfrente os preços $p_1=R\$~2$ por GB e $p_2=R\$~1$ por bloco de minutos. Derivando a demanda marshalliana e avaliando-a nesse ponto, qual é a elasticidade-preço da demanda de $x_1$ em relação ao próprio preço $p_1$?

\begin{enumerate}[label=\Alph*., leftmargin=3.2em, itemsep=0.05em, topsep=0.1em]
  \item $-\dfrac{2}{3}$
  \item $-1$
  \item $0$
  \item $-\dfrac{1}{3}$
  \item $+\dfrac{2}{3}$
\end{enumerate}
\gabaritoitem{A}
\begin{solucao}
\textbf{Passo 1 --- o que se pede.} As preferências são de \emph{complementos perfeitos} (Leontief). Pede-se derivar a demanda marshalliana e calcular a elasticidade-preço própria $\varepsilon_{11}=\dfrac{\partial x_1^{M}}{\partial p_1}\cdot\dfrac{p_1}{x_1^{M}}$ no ponto dado.

\textbf{Passo 2 --- demanda marshalliana.} Com $u=\min\{x_1,x_2\}$, o ótimo nunca desperdiça nenhum bem: consome-se na proporção fixa $x_1=x_2$. Não há condição de tangência (o kink torna a $\text{TMS}$ indefinida no vértice); a otimalidade vem de que comprar mais de um só bem não eleva $u$. Impondo $x_1=x_2$ no orçamento $p_1x_1+p_2x_2=m$:
\[
p_1x_1+p_2x_1=m\;\Rightarrow\; x_1^{M}=\frac{m}{p_1+p_2}.
\]
No ponto: $x_1^{M}=\dfrac{60}{2+1}=20$.

\textbf{Passo 3 --- elasticidade.} Derivando, $\dfrac{\partial x_1^{M}}{\partial p_1}=-\dfrac{m}{(p_1+p_2)^2}$. Logo
\[
\varepsilon_{11}=\frac{\partial x_1^{M}}{\partial p_1}\cdot\frac{p_1}{x_1^{M}}=-\frac{m}{(p_1+p_2)^2}\cdot\frac{p_1}{\,m/(p_1+p_2)\,}=-\frac{p_1}{p_1+p_2}.
\]
O termo $m$ cancela. Com $p_1=2$ e $p_2=1$: $\varepsilon_{11}=-\dfrac{2}{3}$.

\textbf{Passo 4 --- distratores tentadores.} O distrator $-1$ é o erro de quem importa a intuição Cobb-Douglas (elasticidade-preço unitária), esquecendo que em Leontief a presença do bem complementar no denominador $(p_1+p_2)$ amortece a resposta abaixo de $1$ em módulo. O distrator $0$ seduz quem raciocina ``complementos perfeitos não admitem substituição, logo a quantidade não reage ao preço'': confunde ausência de \emph{efeito-substituição} com ausência de resposta --- mas resta o efeito-renda, que faz $x_1$ cair quando $p_1$ sobe. O distrator $-\tfrac13$ troca $p_1$ por $p_2$ no numerador; o $+\tfrac23$ acerta o módulo mas inverte o sinal, violando a lei da demanda.

\textbf{Intuição econômica:} dados e minutos só servem juntos, então o assinante compra ``kits'' $1{:}1$ e gasta $R\$~3$ por kit. Encarecer o GB encarece o kit inteiro, e ele compra menos kits --- a queda é puro efeito-renda, pois não há para onde substituir. Como o GB responde por dois terços do custo do kit, a sensibilidade da quantidade ao seu preço é exatamente $-p_1/(p_1+p_2)=-\tfrac23$: nem tão rígida quanto $0$, nem tão elástica quanto a unitária do caso Cobb-Douglas.
\end{solucao}

\textbf{b)} (3 pontos) O encarecimento do café arábica no varejo (suponha repassado pela Três Corações) leva um analista a modelar a demanda do consumidor a partir da \emph{utilidade indireta} estimada $v(p_1,p_2,m)=\dfrac{m}{p_1^{2/3}\,p_2^{1/3}}$, em que $x_1$ é café (em pacotes) e $x_2$ um bem composto. O expoente $2/3$ sobre $p_1$ reflete uma situação plausível em domicílios de baixa renda, em que o café absorve fração não-trivial do gasto com bebidas. Aplicando a \emph{identidade de Roy} para recuperar a demanda marshalliana por café, e avaliando-a em $m=R\$~90$ e $p_1=R\$~3$, qual é a quantidade demandada de café?

\begin{enumerate}[label=\Alph*., leftmargin=3.2em, itemsep=0.05em, topsep=0.1em]
  \item $30$
  \item $10$
  \item $15$
  \item $20$
  \item $45$
\end{enumerate}
\gabaritoitem{D}
\begin{solucao}
\textbf{Passo 1 --- o que se pede.} A identidade de Roy recupera a demanda marshalliana a partir da utilidade indireta sem reabrir a UMP:
\[
x_1^{M}(p,m)=-\frac{\partial v/\partial p_1}{\partial v/\partial m}.
\]

\textbf{Passo 2 --- derivadas parciais.} De $v=m\,p_1^{-2/3}p_2^{-1/3}$,
\[
\frac{\partial v}{\partial p_1}=m\left(-\tfrac23\right)p_1^{-5/3}p_2^{-1/3},\qquad \frac{\partial v}{\partial m}=p_1^{-2/3}p_2^{-1/3}.
\]

\textbf{Passo 3 --- aplicar Roy.} A razão cancela $p_2^{-1/3}$ e parte de $p_1$:
\[
x_1^{M}=-\frac{m\left(-\tfrac23\right)p_1^{-5/3}p_2^{-1/3}}{p_1^{-2/3}p_2^{-1/3}}=\frac23\,m\,p_1^{-5/3+2/3}=\frac23\cdot\frac{m}{p_1}=\frac{2m}{3p_1}.
\]

\textbf{Passo 4 --- avaliação e distratores.} Em $m=90,\ p_1=3$: $x_1^{M}=\dfrac{2\cdot 90}{3\cdot 3}=\dfrac{180}{9}=20$. O distrator $10$ é o erro de quem usa o expoente $1/3$ de $p_1$ como participação orçamentária (em vez de $2/3$), obtendo $\tfrac13\cdot\tfrac{m}{p_1}=10$; é a troca clássica entre o expoente do bem e o do bem composto. O distrator $30$ resulta de esquecer o fator de participação e reportar $m/p_1=30$ --- o aluno aplica Roy mecanicamente mas perde o coeficiente $\tfrac23$ no cancelamento dos expoentes. O distrator $15$ assume simetria Cobb-Douglas $\tfrac12\cdot\tfrac{m}{p_1}$, e o $45$ esquece o preço no denominador ($m/2$).

\textbf{Intuição econômica:} a utilidade indireta já contém toda a informação do comportamento ótimo --- a identidade de Roy só a ``desempacota''. O expoente $2/3$ sobre $p_1$ sinaliza que o café absorve dois terços do orçamento, o que casa com a ancoragem empírica de domicílios em que a bebida pesa mais no bolso; multiplicado por $m/p_1$, isso dá a quantidade física comprada. É a face dual do problema: em vez de maximizar utilidade sujeito ao orçamento, lemos a demanda diretamente da função-valor, derivando-a em preço e renda.
\end{solucao}

\textbf{c)} (2 pontos, V ou F) Considere a cesta de uma família que aloca todo o orçamento entre vários alimentos (carne bovina, frango, arroz, entre outros) e suponha que a carne bovina seja um \emph{bem de luxo}, isto é, tenha elasticidade-renda maior que $1$ (quando a renda sobe $1\%$, o gasto com carne bovina sobe mais de $1\%$). Sob as convenções do curso (preferências localmente não saciadas, orçamento integralmente exaurido), pode-se afirmar que \emph{necessariamente} pelo menos um outro bem da cesta tem elasticidade-renda menor que $1$.

\gabaritoitem{V}
\begin{solucao}
\textbf{Passo 1 --- a identidade de agregação de Engel.} Diferenciando a restrição orçamentária $\sum_i p_i x_i^{M}(p,m)=m$ em relação a $m$ obtém-se $\sum_i p_i \dfrac{\partial x_i^{M}}{\partial m}=1$. Multiplicando e dividindo cada termo por $x_i^{M}$ e por $m$, isso vira a \emph{agregação de Engel}:
\[
\sum_i s_i\,\eta_i=1,\qquad s_i=\frac{p_i x_i^{M}}{m}\ \ (\text{participações, com }\textstyle\sum_i s_i=1),\quad \eta_i=\frac{\partial x_i^{M}}{\partial m}\frac{m}{x_i^{M}}.
\]

\textbf{Passo 2 --- o argumento.} A média das elasticidades-renda, \emph{ponderada pelas participações}, é exatamente $1$. Se a carne bovina tem $\eta>1$, ela puxa a média ponderada para cima. Para que a média volte a $1$, é impossível que \emph{todos} os demais bens tenham $\eta\ge 1$: nesse caso $\sum_i s_i\eta_i>1$, contradição (com participações positivas). Logo existe ao menos um bem com $\eta_i<1$.

\textbf{Passo 3 --- por que vale ``necessariamente''.} A conclusão não depende da forma funcional nem dos números: é consequência puramente contábil de o orçamento ser exaurido. Basta haver mais de um bem com participação positiva. Note que o teorema garante a \emph{existência} de algum bem com $\eta_i<1$, sem identificá-lo: não se afirma que seja o frango, o arroz ou qualquer item específico. Por isso a afirmação é \textbf{verdadeira}.

\textbf{Intuição econômica:} renda extra precisa ir para algum lugar, e a soma dos gastos sempre fecha em $100\%$. Se a família destina uma fatia \emph{mais que proporcional} do aumento de renda à carne bovina (luxo), sobra menos que proporcional para o conjunto dos outros itens --- algum deles cresce abaixo da renda (necessidade) ou até encolhe (inferior). O frango é o candidato natural: quando a família enriquece, troca parte do frango por carne bovina. Não pode haver economia ``só de luxos''; luxos e necessidades coexistem por força da restrição orçamentária.
\end{solucao}

\textbf{d)} (2 pontos, V ou F) Um consumidor escolhe entre dois botões de gás de cozinha de marcas distintas (suponha Ultragaz, $x_1$, e uma concorrente, $x_2$), que ele considera \emph{substitutos perfeitos} segundo $u(x_1,x_2)=2x_1+3x_2$. Suponha que os preços sejam $p_1=R\$~4$ e $p_2=R\$~3$ e que a renda destinada a gás seja $m>0$. Sob as convenções do curso, a cesta ótima envolve quantidades estritamente positivas de \emph{ambas} as marcas, pois diversificar entre fornecedores aumenta a utilidade.

\gabaritoitem{F}
\begin{solucao}
\textbf{Passo 1 --- a regra dos substitutos perfeitos.} Com $u=2x_1+3x_2$, a $\text{TMS}$ é constante: $\text{TMS}=\dfrac{\partial u/\partial x_1}{\partial u/\partial x_2}=\dfrac{2}{3}$. As curvas de indiferença são retas paralelas. A decisão se resume a comparar a \emph{utilidade por real} de cada marca.

\textbf{Passo 2 --- utilidade por real.} 
\[
\frac{\partial u/\partial x_1}{p_1}=\frac{2}{4}=0{,}5,\qquad \frac{\partial u/\partial x_2}{p_2}=\frac{3}{3}=1{,}0.
\]
A marca $2$ entrega o dobro de utilidade por real gasto. O consumidor racional concentra \emph{todo} o orçamento nela: $x_2^{*}=m/3$ e $x_1^{*}=0$. Solução de \emph{canto}.

\textbf{Passo 3 --- por que a afirmação falha.} Como $\text{TMS}=\tfrac23\neq \tfrac{p_1}{p_2}=\tfrac43$, não há tangência interior; o ótimo é de canto, não interior. Comprar a marca $1$ desperdiçaria orçamento. A justificativa de ``diversificar fornecedores'' é economicamente espelho de uma curva de indiferença \emph{estritamente convexa} (gosto por mistura), que \emph{não existe} em substitutos perfeitos: aqui as indiferenças são lineares, sem nenhuma vantagem em misturar. Logo a afirmação é \textbf{falsa}.

\textbf{Intuição econômica:} substitutos perfeitos são como trocar moedas a uma taxa fixa. Se uma marca rende mais utilidade por real, o consumidor vai inteiramente para ela --- não existe motivo de variedade quando os bens são intercambiáveis a taxa constante. A única exceção seria o empate $\text{TMS}=p_1/p_2$, caso em que qualquer divisão é ótima (e mesmo aí o interior não é \emph{necessário}). Com os preços dados, o empate não ocorre, e a solução de canto é única.
\end{solucao}

\ifsolucao\else
\par\vspace{0.5em}
\begin{center}
\textbf{Quadro-resposta do Bloco 1 --- escreva a alternativa (A--E) ou V/F de cada item:}\\[0.6em]
\renewcommand{\arraystretch}{1.45}
\begin{tabular}{|>{\centering\arraybackslash}p{2.6cm}|>{\centering\arraybackslash}p{1.4cm}|>{\centering\arraybackslash}p{1.4cm}|>{\centering\arraybackslash}p{1.4cm}|>{\centering\arraybackslash}p{1.4cm}|}
  \hline
  \textbf{Item} & a) & b) & c) & d) \\ \hline
  \textbf{Resposta} &  &  &  &  \\ \hline
\end{tabular}
\end{center}
{\small\itshape Atenção: apenas o quadro acima será considerado na correção dos itens objetivos deste bloco.}\par\vspace{0.6em}
\fi

\textbf{Questão 2.} (15 pontos) Com a popularização do delivery, considere um consumidor que reparte seu orçamento mensal de jantares entre \emph{refeições pedidas por aplicativo} (suponha o iFood), $x_1$, e \emph{refeições preparadas em casa}, $x_2$. Os dois bens são substitutos \emph{imperfeitos}: trocar app por cozinha caseira é possível, mas há conveniência que não se transfere integralmente. Modele as preferências pela função CES simétrica $u(x_1,x_2)=\left(x_1^{\rho}+x_2^{\rho}\right)^{1/\rho}$, com $\rho=\tfrac12$. Suponha que o consumidor disponha de $m=R\$~48$ no mês e enfrente os preços $p_1=R\$~2$ por refeição de app e $p_2=R\$~4$ por refeição caseira (em ``unidades'' comparáveis de uma refeição). Em todos os itens, considere solução interior.

\textbf{a)} (5 pontos) Monte a UMP e derive as demandas marshallianas $x_1^{M}(p,m)$ e $x_2^{M}(p,m)$ em forma fechada para a CES simétrica com $\rho=\tfrac12$ (equivalentemente, elasticidade de substituição $\sigma=2$). Em seguida, avalie-as nos preços e renda dados e verifique que a restrição orçamentária é satisfeita.
\resolucaobox{5.5cm}
\begin{solucao}
\textbf{Passo 1 --- UMP e CPO.} O consumidor resolve $\max\ (x_1^{\rho}+x_2^{\rho})^{1/\rho}$ s.a. $p_1x_1+p_2x_2=m$. Igualando a $\text{TMS}$ à razão de preços,
\[
\text{TMS}=\frac{x_1^{\rho-1}}{x_2^{\rho-1}}=\left(\frac{x_1}{x_2}\right)^{\rho-1}=\frac{p_1}{p_2}\;\Rightarrow\;\frac{x_1}{x_2}=\left(\frac{p_1}{p_2}\right)^{1/(\rho-1)}=\left(\frac{p_1}{p_2}\right)^{-\sigma},
\]
com $\sigma=\dfrac{1}{1-\rho}=\dfrac{1}{1-1/2}=2$.

\textbf{Passo 2 --- forma fechada.} Levando $x_2=x_1(p_1/p_2)^{\sigma}$ ao orçamento e isolando,
\[
x_1^{M}=\frac{m\,p_1^{-\sigma}}{p_1^{\,1-\sigma}+p_2^{\,1-\sigma}},\qquad x_2^{M}=\frac{m\,p_2^{-\sigma}}{p_1^{\,1-\sigma}+p_2^{\,1-\sigma}}.
\]
Com $\sigma=2$ temos $1-\sigma=-1$, logo o denominador é $p_1^{-1}+p_2^{-1}$.

\textbf{Passo 3 --- avaliação.} Em $p_1=2,\ p_2=4,\ m=48$: denominador $=\tfrac12+\tfrac14=\tfrac34$. Então
\[
x_1^{M}=\frac{48\cdot 2^{-2}}{3/4}=\frac{48\cdot \tfrac14}{3/4}=\frac{12}{3/4}=16,\qquad x_2^{M}=\frac{48\cdot 4^{-2}}{3/4}=\frac{48\cdot \tfrac1{16}}{3/4}=\frac{3}{3/4}=4.
\]

\textbf{Passo 4 --- verificação orçamentária.} $p_1x_1^{M}+p_2x_2^{M}=2\cdot16+4\cdot4=32+16=48=m.$ Confere, e ambas as quantidades são positivas (solução interior).

\textbf{Intuição econômica:} a CES interpola entre rigidez e flexibilidade de escolha. Como o app está mais barato por refeição ($R\$~2$ contra $R\$~4$), o consumidor pende para ele ($16$ contra $4$), mas não vai ao canto: a imperfeição da substituição ($\sigma=2$ finito) preserva alguma cozinha caseira na cesta. O grau dessa inclinação é governado por $\sigma$.
\rubricalabel
\begin{itemize}
\item Monta a UMP corretamente (objetivo CES $+$ restrição): 1 pt.
\item Obtém a razão ótima $x_1/x_2=(p_1/p_2)^{-\sigma}$ a partir da $\text{TMS}$, com $\sigma=1/(1-\rho)=2$: 1{,}5 pt.
\item Escreve a forma fechada $x_i^{M}=m\,p_i^{-\sigma}/(p_1^{1-\sigma}+p_2^{1-\sigma})$: 1 pt.
\item Avalia $x_1^{M}=16$ e $x_2^{M}=4$: 1 pt (0{,}5 cada).
\item Verifica que o orçamento fecha em $48$ (interioridade): 0{,}5 pt.
\item Erro aritmético isolado com método integralmente correto: penalidade leve de $-0{,}5$ pt (não zera o item).
\end{itemize}
\end{solucao}

\textbf{b)} (5 pontos) Mostre que, para esta CES, a \emph{elasticidade de substituição} entre os dois bens vale $\sigma=2$ e interprete o número. Em seguida, usando que a participação orçamentária do app no ponto ótimo é $s_1=p_1x_1^{M}/m$, calcule a \emph{elasticidade-preço própria} da demanda por refeições de app, $\varepsilon_{11}$, e classifique a demanda como elástica ou inelástica.
\resolucaobox{5.5cm}
\begin{solucao}
\textbf{Passo 1 --- elasticidade de substituição.} Por definição, $\sigma=-\dfrac{d\ln(x_1/x_2)}{d\ln(\text{TMS})}$. Do item anterior, $\dfrac{x_1}{x_2}=\left(\dfrac{p_1}{p_2}\right)^{-\sigma}$ e, no ótimo, $\text{TMS}=p_1/p_2$. Tomando logaritmos, $\ln(x_1/x_2)=-\sigma\ln(\text{TMS})$, donde a inclinação é $-\sigma$ e portanto $\sigma=\dfrac{1}{1-\rho}=\dfrac{1}{1-1/2}=2$.

\textbf{Passo 2 --- interpretação.} $\sigma=2$ significa que uma variação de $1\%$ na razão de preços $p_1/p_2$ altera a razão de quantidades $x_1/x_2$ em $2\%$ (sentido oposto). É substituibilidade \emph{moderada e finita}: maior que a Cobb-Douglas ($\sigma=1$) e os complementos ($\sigma=0$), menor que substitutos perfeitos ($\sigma\to\infty$).

\textbf{Passo 3 --- participação e elasticidade própria.} A participação do app é $s_1=\dfrac{p_1x_1^{M}}{m}=\dfrac{2\cdot16}{48}=\dfrac{32}{48}=\dfrac23$. Para a CES (homotética, $\eta_1=1$), a elasticidade-preço própria é
\[
\varepsilon_{11}=-\sigma(1-s_1)-s_1=-\sigma+(\sigma-1)s_1.
\]
Com $\sigma=2$ e $s_1=\tfrac23$: $\varepsilon_{11}=-2+(2-1)\cdot\tfrac23=-2+\tfrac23=-\dfrac{4}{3}.$

\textbf{Passo 4 --- classificação.} Como $|\varepsilon_{11}|=\tfrac43>1$, a demanda por refeições de app é \emph{elástica} no ponto considerado.

\textbf{Intuição econômica:} a elasticidade própria decompõe-se em duas forças. A substituição pura empurra com peso $-\sigma(1-s_1)$: quanto mais fácil migrar para a cozinha caseira (maior $\sigma$) e quanto menor a fatia do app, mais sensível a demanda. O termo $-s_1$ é o efeito-renda. Com $\sigma=2$ e o app já ocupando dois terços do orçamento, o resultado líquido é $-\tfrac43$: encarecer o app em $1\%$ derruba a quantidade pedida em cerca de $1{,}33\%$.
\rubricalabel
\begin{itemize}
\item Mostra $\sigma=1/(1-\rho)=2$ via $\ln(x_1/x_2)=-\sigma\ln(\text{TMS})$ (ou definição equivalente): 1{,}5 pt.
\item Interpreta corretamente o significado de $\sigma=2$ (resposta proporcional da razão de quantidades): 1 pt.
\item Calcula a participação $s_1=2/3$: 0{,}5 pt.
\item Obtém $\varepsilon_{11}=-\sigma+(\sigma-1)s_1=-4/3$ (aceita-se dedução direta por derivação da demanda fechada): 1{,}5 pt.
\item Classifica como demanda elástica ($|\varepsilon_{11}|>1$): 0{,}5 pt.
\item Erro aritmético isolado com método correto (ex.: fórmula certa, conta final trocada): penalidade leve de $-0{,}5$ pt.
\item Usar fórmula sem efeito-renda ($\varepsilon_{11}=-\sigma$) sem justificar a omissão: teto de 3 pts no item.
\end{itemize}
\end{solucao}

\textbf{c)} (5 pontos) (i) Prove que esta CES é \emph{homotética}, mostrando que as demandas são lineares na renda (equivalentemente, $x_i^{M}(p,tm)=t\,x_i^{M}(p,m)$) e conclua que a elasticidade-renda é $\eta_i=1$. (ii) Indique para que formas funcionais a CES degenera nos limites $\rho\to 1$, $\rho\to 0$ e $\rho\to -\infty$, dando o valor de $\sigma$ em cada caso (basta nomear a forma e o $\sigma$; uma breve justificativa por caso é suficiente). (iii) Para o limite $\rho\to 0$, mostre \emph{rigorosamente} que $u$ converge à Cobb-Douglas $x_1^{1/2}x_2^{1/2}$ (a menos de transformação monótona). Esta é a única demonstração formal exigida no item.
\resolucaobox{5.5cm}
\begin{solucao}
\textbf{Passo 1 --- homoteticidade.} Na forma fechada $x_i^{M}=\dfrac{m\,p_i^{-\sigma}}{p_1^{1-\sigma}+p_2^{1-\sigma}}$, a renda $m$ entra como \emph{fator multiplicativo}: $x_i^{M}(p,tm)=\dfrac{tm\,p_i^{-\sigma}}{p_1^{1-\sigma}+p_2^{1-\sigma}}=t\,x_i^{M}(p,m)$. Logo a demanda é linear em $m$ a preços fixos, as curvas de Engel são retas pela origem, e
\[
\eta_i=\frac{\partial x_i^{M}}{\partial m}\frac{m}{x_i^{M}}=\frac{x_i^{M}/m\cdot m}{x_i^{M}}=1.
\]
Equivalentemente, $u$ é homogênea de grau $1$, logo homotética, e preferências homotéticas têm expansão radial: dobrar a renda dobra cada quantidade.

\textbf{Passo 2 --- limites de $\sigma$.} Como $\sigma=\dfrac{1}{1-\rho}$:
\begin{itemize}
\item $\rho\to 1\Rightarrow \sigma\to\infty$: $u\to x_1+x_2$, \emph{substitutos perfeitos} (curvas de indiferença retas). Justificativa: $\rho=1$ dá diretamente a soma.
\item $\rho\to 0\Rightarrow \sigma\to 1$: \emph{Cobb-Douglas} (demonstração rigorosa no Passo 3).
\item $\rho\to -\infty\Rightarrow \sigma\to 0$: $u\to\min\{x_1,x_2\}$, \emph{complementos perfeitos} (Leontief). Justificativa: para $\rho\to-\infty$, em $(x_1^{\rho}+x_2^{\rho})^{1/\rho}$ o termo de \emph{maior} expoente negativo domina, ou seja, o da \emph{menor} quantidade, e o limite é o mínimo.
\end{itemize}

\textbf{Passo 3 --- limite Cobb-Douglas ($\rho\to 0$).} Considere $\ln u=\dfrac{1}{\rho}\ln(x_1^{\rho}+x_2^{\rho})$. Quando $\rho\to 0$, $\ln(x_1^{\rho}+x_2^{\rho})\to\ln(x_1^{0}+x_2^{0})=\ln 2$, de modo que $\ln u\to\dfrac{\ln 2}{0}$, isto é, \emph{diverge} --- não é indeterminação $0/0$. Subtraio então a parcela $\dfrac{\ln 2}{\rho}$, que é \emph{constante em} $(x_1,x_2)$, formando a diferença
\[
g(\rho)=\frac{\ln(x_1^{\rho}+x_2^{\rho})-\ln 2}{\rho},
\]
cujo numerador $\to \ln 2-\ln 2=0$ e cujo denominador $\to 0$: \emph{agora sim} a forma é $0/0$, à qual aplico L'Hôpital. Subtrair $\tfrac{\ln 2}{\rho}$ é legítimo porque essa parcela não depende de $(x_1,x_2)$ e apenas desloca aditivamente $\ln u$, preservando a ordenação de preferências. Derivando numerador e denominador em $\rho$, com $\dfrac{d}{d\rho}x_i^{\rho}=x_i^{\rho}\ln x_i$:
\[
\lim_{\rho\to 0}g(\rho)=\lim_{\rho\to 0}\frac{x_1^{\rho}\ln x_1+x_2^{\rho}\ln x_2}{x_1^{\rho}+x_2^{\rho}}=\frac{\ln x_1+\ln x_2}{2}=\ln\!\big(x_1^{1/2}x_2^{1/2}\big).
\]
Logo, a menos da parcela constante (em $x$) $\tfrac{\ln 2}{\rho}$, $\ln u\to \ln(x_1^{1/2}x_2^{1/2})$, ou seja, $u\to x_1^{1/2}x_2^{1/2}$. A parcela $\tfrac{\ln 2}{\rho}$ é uma constante que não depende de $(x_1,x_2)$; somá-la ou subtraí-la equivale a uma transformação monótona (adição de constante) da função de utilidade, que preserva a ordenação de preferências. Como a CES simétrica reparte o orçamento igualmente no limite, recupera-se a Cobb-Douglas de expoentes $\tfrac12,\tfrac12$.

\textbf{Intuição econômica:} a CES é a ``família-mãe'' das formas funcionais do curso: um único parâmetro $\rho$ (ou $\sigma$) percorre todo o espectro de substituibilidade. O delivery e a cozinha caseira aparecem como casos intermediários entre o cenário dos planos de celular (Leontief, $\sigma=0$, este bloco), o do café modelado em Cobb-Douglas ($\sigma=1$) e o do gás de cozinha em substitutos perfeitos ($\sigma\to\infty$). A homoteticidade, presente em toda a família CES, garante que mudanças de renda deslocam a cesta radialmente, sem alterar as participações orçamentárias --- por isso $\eta_i=1$ para todos os bens.
\rubricalabel
\begin{itemize}
\item Prova homoteticidade via linearidade em $m$ ($x_i(p,tm)=t\,x_i(p,m)$) e conclui $\eta_i=1$: 1{,}5 pt.
\item Identifica os três limites com $\sigma$ correto ($\rho\to1$: subst. perfeitos $\sigma\to\infty$; $\rho\to0$: Cobb-Douglas $\sigma=1$; $\rho\to-\infty$: Leontief $\sigma=0$): 1{,}5 pt (0{,}5 cada). Para os limites $\rho\to1$ e $\rho\to-\infty$ basta nomear forma e $\sigma$ com breve justificativa.
\item Demonstra rigorosamente o limite Cobb-Douglas, montando corretamente a forma $0/0$ (subtração da constante $\ln2/\rho$) e aplicando L'Hôpital (ou expansão equivalente sobre $\ln u$): 1{,}5 pt.
\item Observa que a parcela $\ln2/\rho$ é constante em $(x_1,x_2)$ e que somá-la/subtraí-la é transformação monótona que preserva a ordenação: 0{,}5 pt.
\item Apenas \emph{afirmar} o limite Cobb-Douglas sem montar a indeterminação nem justificar: teto de 2 pts no item.
\item Erro aritmético/de sinal isolado em derivada com método correto na demonstração do limite: penalidade leve de $-0{,}5$ pt.
\end{itemize}
\end{solucao}


\blocotitulo{Bloco 2 --- Equilíbrio Geral \hfill \normalsize (25 pontos)}

\textbf{Questão 3.} (10 pontos) \textbf{Itens objetivos.} Nos itens \textbf{a)} e \textbf{b)}, assinale a \emph{única} alternativa correta (3 pontos cada; não há penalização por erro). Nos itens \textbf{c)} e \textbf{d)}, responda \textbf{Verdadeiro ou Falso} (2 pontos cada; vale a regra de anulação descrita na capa). Registre suas respostas no quadro-resposta ao final do bloco.

\textbf{a)} (3 pontos) A Vale aloca eficientemente um insumo escasso entre minério de ferro para exportação ($x_1$) e pelotas (minério pelotizado, $x_2$), enfrentando a fronteira de possibilidades de produção (FPP) $x_1^2 + x_2^2 = 25$, com $x_1,x_2\ge 0$. Aos preços de mercado $p_1 = 3$ e $p_2 = 4$, o plano de produção sobre a FPP que maximiza a receita $p_1 x_1 + p_2 x_2$ é:

\begin{enumerate}[label=\Alph*., leftmargin=3.2em, itemsep=0.05em, topsep=0.1em]
  \item $(x_1^{*},x_2^{*}) = (5,0)$
  \item $(x_1^{*},x_2^{*}) = (4,3)$
  \item $(x_1^{*},x_2^{*}) = (3,4)$
  \item $(x_1^{*},x_2^{*}) = (0,5)$
  \item $(x_1^{*},x_2^{*}) = \big(\sqrt{12{,}5},\,\sqrt{12{,}5}\big)$
\end{enumerate}
\gabaritoitem{C}
\begin{solucao}
\textbf{Passo 1 --- O que é dado e o que se pede.} FPP côncava $x_1^2+x_2^2=25$ (um quarto de círculo de raio $5$), preços $p_1=3$, $p_2=4$. Pede-se o plano que maximiza a receita $p_1 x_1 + p_2 x_2$ sobre a FPP.

\textbf{Passo 2 --- Plano.} O ótimo de produção iguala a $\text{TMT}$ à razão de preços. Ao longo da FPP, $x_2 = \sqrt{25 - x_1^2}$, logo $\dfrac{dx_2}{dx_1} = -\dfrac{x_1}{\sqrt{25-x_1^2}} = -\dfrac{x_1}{x_2}$ e portanto $\text{TMT} = \dfrac{x_1}{x_2}$. A condição de receita máxima é $\text{TMT}=p_1/p_2$.

\textbf{Passo 3 --- Execução.} Igualando $\text{TMT}=p_1/p_2$:
\[ \frac{x_1}{x_2} = \frac{3}{4} \;\Longrightarrow\; x_2 = \frac{4}{3}x_1. \]
Substituindo na FPP: $x_1^2 + \dfrac{16}{9}x_1^2 = 25 \Rightarrow \dfrac{25}{9}x_1^2 = 25 \Rightarrow x_1^2 = 9 \Rightarrow x_1^{*} = 3$ e $x_2^{*} = 4$. (Triângulo $3$--$4$--$5$.) Plano eficiente $(3,4)$ --- alternativa (C).

\textbf{Passo 4 --- Verificação.} A receita em $(3,4)$ é $3\cdot 3 + 4\cdot 4 = 9 + 16 = 25$. Compare: $(4,3)$ rende $12+12 = 24$; $(5,0)$ rende $15$; $(0,5)$ rende $20$. Nenhum ponto da FPP supera $25$, confirmando o ótimo.

\textbf{Por que os distratores mais tentadores enganam.} A alternativa (B), $(4,3)$, é a mais sedutora: resulta de \emph{inverter} a razão de preços, igualando $\text{TMT}=p_2/p_1=4/3$ em vez de $p_1/p_2=3/4$ --- erro de quem confunde qual preço vai no numerador da $\text{TMT}$. A alternativa (E), $\big(\sqrt{12{,}5},\sqrt{12{,}5}\big)$, vem de \emph{dividir o insumo igualmente} entre as duas linhas (ponto médio da FPP), ignorando que preços distintos pedem produção desigual: ali $\text{TMT}=1\neq 3/4$, e há receita não explorada deslocando produção para a pelota, mais valiosa. As alternativas (A) e (D) são as especializações de canto, que desperdiçam o bem omitido apesar de ele ter preço positivo.

\textbf{Intuição econômica:} a curvatura da FPP reflete custo de oportunidade crescente --- cada tonelada extra de minério para exportação custa, na margem, cada vez mais pelotas sacrificadas. O plano eficiente para a Vale é aquele em que a taxa \emph{técnica} de transformação entre os bens iguala a taxa \emph{econômica} a que o mercado os troca; produzir mais pelotas, relativamente mais caras, é racional até que a tecnologia torne essa conversão cara demais.
\end{solucao}

\textbf{b)} (3 pontos) No entreposto da CEAGESP, em São Paulo, produtores trocam livremente dois bens hortícolas a preços tomados como dados: tomate ($x_1$) e uma cesta composta de hortaliças ($x_2$). Modelado como economia de trocas, o excesso de demanda agregado pelo bem $1$ é $z_1(p_1,p_2) = \dfrac{600\,p_2}{p_1} - 200$, dependente apenas da razão de preços. Tomando a cesta composta como numerário ($p_2 = 1$), considere as afirmações sobre o equilíbrio walrasiano desse entreposto. Assinale a \emph{correta}.

\begin{enumerate}[label=\Alph*., leftmargin=3.2em, itemsep=0.05em, topsep=0.1em]
  \item O único equilíbrio é o vetor absoluto $(p_1,p_2) = (3,1)$; qualquer outro vetor de preços viola o equilíbrio.
  \item Dobrar todos os preços para $(p_1,p_2)=(6,2)$ dobra o excesso de demanda, pois $z_1$ é homogênea de grau um.
  \item O preço relativo de equilíbrio é $p_1/p_2 = 1/3$, de modo que o tomate vale um terço de uma unidade da cesta composta.
  \item A escolha do bem numerário altera a alocação real de equilíbrio, pois fixar $p_2=1$ injeta valor real na economia.
  \item O equilíbrio é determinado apenas pela razão $p_1/p_2 = 3$; o vetor $(6,2)$ descreve \emph{o mesmo} equilíbrio que $(3,1)$, pois $z_1$ é homogênea de grau zero.
\end{enumerate}
\gabaritoitem{E}
\begin{solucao}
\textbf{Passo 1 --- O que é dado e o que se pede.} Excesso de demanda $z_1(p)=\dfrac{600 p_2}{p_1}-200$, função apenas da razão $p_2/p_1$. Pede-se a afirmação correta sobre o equilíbrio e o papel da normalização.

\textbf{Passo 2 --- Equilíbrio.} Equilíbrio exige $z_1(p)=0$:
\[ \frac{600 p_2}{p_1} = 200 \;\Longrightarrow\; \frac{p_1}{p_2} = 3. \]
Apenas a \emph{razão} é determinada. Com $p_2=1$ (numerário), $p_1=3$.

\textbf{Passo 3 --- Homogeneidade de grau zero.} Como $z_1$ depende só de $p_2/p_1$, para qualquer $\lambda>0$ vale $z_1(\lambda p_1,\lambda p_2)=z_1(p_1,p_2)$: $z$ é homogênea de \emph{grau zero}. Logo $(p_1,p_2)=(3,1)$ e $(6,2)$ representam o \emph{mesmo} equilíbrio --- mesma razão $3$, mesmas demandas, mesma alocação real. A afirmação (E) é a correta.

\textbf{Passo 4 --- Verificação.} Em $(6,2)$: $z_1 = \dfrac{600\cdot 2}{6} - 200 = 200 - 200 = 0$. Confirma o equilíbrio sem qualquer mudança real; a normalização é mera escolha de unidade de conta.

\textbf{Por que os distratores mais tentadores enganam.} A alternativa (A) é a mais atraente: confunde a \emph{normalização} (uma convenção que fixa um preço) com \emph{unicidade do vetor absoluto}. Mas só razões de preços têm conteúdo econômico; escalar todos os preços não cria novo equilíbrio. A alternativa (B) troca grau zero por grau um --- erro de quem lembra a palavra ``homogênea'' mas erra o grau; se fosse grau um, dobrar preços dobraria $z$, o que aqui é falso. A alternativa (C) inverte a razão ($1/3$ em vez de $3$). A alternativa (D) confunde fixar o numerário com alterar dotações reais.

\textbf{Intuição econômica:} no entreposto, só preços \emph{relativos} importam --- o dinheiro é véu. A homogeneidade de grau zero do excesso de demanda é a expressão formal disso: multiplicar todos os preços por um fator comum não muda nem orçamentos relativos, nem demandas, nem o equilíbrio. Escolher um numerário (cesta composta $=1$) apenas ancora a régua; a troca real entre os produtores é idêntica meça-se em cestas ou em qualquer outra unidade.
\end{solucao}

\textbf{c)} (2 pontos, V ou F) Em uma região cafeeira de Minas Gerais, cooperativas de produtores trocam entre si dois bens --- café arábica ($x_1$) e um bem composto ($x_2$) --- em uma economia de trocas pura. Suponha que as preferências de cada cooperativa sejam contínuas e localmente não saciadas (mas \emph{não} necessariamente convexas) e que todas tomem os preços como dados. \textbf{Afirmação:} para que \emph{todo} equilíbrio walrasiano dessa economia seja Pareto-eficiente --- conteúdo do 1.\textsuperscript{o} Teorema do Bem-Estar --- é \emph{indispensável} supor, adicionalmente, que as preferências de cada cooperativa sejam convexas.

\gabaritoitem{F}
\begin{solucao}
\textbf{Passo 1 --- O que se afirma.} Que o 1.\textsuperscript{o} Teorema do Bem-Estar \emph{exige} convexidade das preferências.

\textbf{Passo 2 --- Enunciado correto.} O 1.\textsuperscript{o} Teorema requer essencialmente uma única hipótese sobre preferências: \emph{não-saciedade local}. Sob ela, todo equilíbrio walrasiano (com agentes tomadores de preço e mercados competitivos) é Pareto-eficiente --- \emph{sem} qualquer suposição de convexidade.

\textbf{Passo 3 --- Por que dispensa convexidade.} A prova é um argumento de revelação de preferências: se uma alocação viável Pareto-dominasse o equilíbrio, então, por não-saciedade local, a cesta preferida de cada agente custaria \emph{pelo menos} sua renda (e estritamente mais para quem melhora), de modo que a alocação dominante custaria, no agregado, mais que a renda total --- contradição com a viabilidade ($\sum_i p\cdot x_i = \sum_i p\cdot\omega_i$). Nenhum passo usa a forma das curvas de indiferença; convexidade nunca entra.

\textbf{Passo 4 --- Onde a convexidade realmente importa.} A convexidade é hipótese essencial do \emph{2.\textsuperscript{o}} Teorema (descentralizar uma alocação eficiente como equilíbrio com transferências), pois é ela que viabiliza a separação por hiperplano. Confundir os dois teoremas é o erro que a afirmação comete --- por isso é \textbf{falsa}.

\textbf{Intuição econômica:} o 1.\textsuperscript{o} Teorema é o lado \emph{robusto} da mão invisível: se as cooperativas apenas ``sempre querem um pouco mais'' (não-saciedade local) e respondem a preços, a barganha competitiva esgota todos os ganhos de troca, chegando à curva de contrato --- mesmo que alguma preferência tenha curvas de indiferença não convexas. A convexidade é o pedágio do caminho \emph{inverso} (Pareto $\Rightarrow$ descentralizável), não deste.
\end{solucao}

\textbf{d)} (2 pontos, V ou F) O comércio no âmbito do Mercosul entre Brasil e Argentina é descrito como troca de trigo ($x_1$, em que a Argentina é relativamente abundante) por bens industrializados ($x_2$). Modele-o como uma economia de trocas $2\times 2$ (agentes $A$ = Brasil e $B$ = Argentina) cujas preferências são contínuas, estritamente monótonas, estritamente convexas e diferenciáveis. \textbf{Afirmação:} uma alocação \emph{interior} dessa caixa de Edgeworth é Pareto-eficiente se, e somente se, as taxas marginais de substituição dos dois países se igualam, $\text{TMS}^A = \text{TMS}^B$; o lugar geométrico de todas essas alocações eficientes é exatamente a curva de contrato.

\gabaritoitem{V}
\begin{solucao}
\textbf{Passo 1 --- O que se afirma.} Que, no interior da caixa, eficiência de Pareto equivale a tangência das curvas de indiferença ($\text{TMS}^A=\text{TMS}^B$), e que o conjunto desses pontos é a curva de contrato.

\textbf{Passo 2 --- Necessidade ($\Rightarrow$).} Suponha eficiência mas $\text{TMS}^A\neq\text{TMS}^B$ num ponto interior. Então existe uma direção de troca marginal $(d x_1, d x_2)$ que aumenta $u^A$ sem reduzir $u^B$ (uma das partes valoriza o trigo mais que a outra na margem, abrindo ganho de troca). Isso contradiz a eficiência. Logo eficiência interior $\Rightarrow \text{TMS}^A=\text{TMS}^B$.

\textbf{Passo 3 --- Suficiência ($\Leftarrow$).} Com preferências estritamente convexas e monótonas, a igualdade das $\text{TMS}$ no interior é também \emph{suficiente}: as curvas de indiferença são tangentes e cada conjunto ``ao menos tão bom'' é convexo, de modo que não há cesta viável que melhore um país sem piorar o outro. A condição de tangência é, portanto, caracterização completa no interior.

\textbf{Passo 4 --- Curva de contrato.} O conjunto de todas as alocações viáveis com $\text{TMS}^A=\text{TMS}^B$ é, por definição, a curva de contrato dentro da caixa de Edgeworth. Daí a afirmação ser \textbf{verdadeira}. (A convexidade estrita é o que garante que a condição de primeira ordem caracterize o ótimo, e não apenas um ponto crítico --- por isso o enunciado a inclui.)

\textbf{Intuição econômica:} enquanto Brasil e Argentina avaliam o trigo a taxas marginais distintas, há um negócio mutuamente vantajoso a fazer --- alguém troca trigo por industrializados e ambos ganham. Os ganhos de troca só se esgotam quando as avaliações marginais coincidem; esse alinhamento das $\text{TMS}$, varrido por todas as divisões possíveis do excedente, desenha a curva de contrato, o conjunto dos acordos eficientes do bloco.
\end{solucao}

\ifsolucao\else
\par\vspace{0.5em}
\begin{center}
\textbf{Quadro-resposta do Bloco 2 --- escreva a alternativa (A--E) ou V/F de cada item:}\\[0.6em]
\renewcommand{\arraystretch}{1.45}
\begin{tabular}{|>{\centering\arraybackslash}p{2.6cm}|>{\centering\arraybackslash}p{1.4cm}|>{\centering\arraybackslash}p{1.4cm}|>{\centering\arraybackslash}p{1.4cm}|>{\centering\arraybackslash}p{1.4cm}|}
  \hline
  \textbf{Item} & a) & b) & c) & d) \\ \hline
  \textbf{Resposta} &  &  &  &  \\ \hline
\end{tabular}
\end{center}
{\small\itshape Atenção: apenas o quadro acima será considerado na correção dos itens objetivos deste bloco.}\par\vspace{0.6em}
\fi

\textbf{Questão 4.} (15 pontos) A Autoridade Portuária de Santos precisa dividir dois insumos escassos e fixos entre dois operadores de terminal, $A$ e $B$: janelas de atracação (bem $1$, em horas-berço) e área de pátio para contêineres (bem $2$, em lotes de armazenagem). A dotação total da semana é $\bar\omega_1 = 12$ horas-berço e $\bar\omega_2 = 12$ lotes de pátio, de modo que a alocação dos operadores pode ser representada em uma caixa de Edgeworth $12\times 12$. As preferências de cada operador sobre os dois insumos são Cobb-Douglas: $u^A(x_1,x_2) = x_1^{2/3}\,x_2^{1/3}$ (o operador $A$ é intensivo em atracação) e $u^B(x_1,x_2) = x_1^{1/3}\,x_2^{2/3}$ (o operador $B$ é intensivo em pátio). Suponha que os operadores possam renegociar livremente o uso dos insumos a preços tomados como dados, com o bem $2$ servindo de numerário ($p_2 = 1$).

\textbf{a)} (5 pontos) Derive \emph{explicitamente} a equação da curva de contrato dessa caixa de Edgeworth, expressando a alocação do operador $A$ na forma $x_2^A = g(x_1^A)$. Use a condição de eficiência $\text{TMS}^A = \text{TMS}^B$ junto com a viabilidade. Verifique que a curva obtida passa pelos dois cantos da caixa (origem de $A$ e origem de $B$).
\resolucaobox{5.5cm}
\begin{solucao}
\textbf{Passo 1 --- Taxas marginais de substituição.} Para $u=x_1^{a}x_2^{1-a}$, $\text{TMS}=\dfrac{a}{1-a}\dfrac{x_2}{x_1}$. Assim,
\[ \text{TMS}^A = \frac{2/3}{1/3}\frac{x_2^A}{x_1^A} = 2\,\frac{x_2^A}{x_1^A}, \qquad \text{TMS}^B = \frac{1/3}{2/3}\frac{x_2^B}{x_1^B} = \frac{1}{2}\,\frac{x_2^B}{x_1^B}. \]

\textbf{Passo 2 --- Viabilidade.} Na caixa $12\times 12$, $x_1^B = 12 - x_1^A$ e $x_2^B = 12 - x_2^A$.

\textbf{Passo 3 --- Igualar as $\text{TMS}$.} A eficiência interior exige $\text{TMS}^A=\text{TMS}^B$:
\[ 2\,\frac{x_2^A}{x_1^A} = \frac{1}{2}\,\frac{12 - x_2^A}{12 - x_1^A}. \]
Cruzando os termos: $4\,x_2^A\,(12 - x_1^A) = x_1^A\,(12 - x_2^A)$, isto é, $48\,x_2^A - 4\,x_1^A x_2^A = 12\,x_1^A - x_1^A x_2^A$. Reagrupando, $48\,x_2^A - 3\,x_1^A x_2^A = 12\,x_1^A$, donde $x_2^A\,(48 - 3\,x_1^A) = 12\,x_1^A$ e
\[ \boxed{\,x_2^A = \frac{12\,x_1^A}{48 - 3\,x_1^A} = \frac{4\,x_1^A}{16 - x_1^A}\,}, \qquad x_1^A\in(0,12). \]

\textbf{Passo 4 --- Passagem pelos cantos.} Em $x_1^A\to 0$: $x_2^A\to 0$ --- origem de $A$ (operador $B$ fica com tudo). Em $x_1^A = 12$: $x_2^A = \dfrac{48}{16-12} = \dfrac{48}{4} = 12$ --- operador $A$ fica com toda a caixa (origem de $B$). A curva é crescente e abaulada, ligando os dois cantos, como esperado para a curva de contrato.

\textbf{Intuição econômica:} a curva de contrato reúne todas as divisões dos insumos do porto em que nenhum operador consegue melhorar sem prejudicar o outro. Como $A$ valoriza relativamente mais as janelas de atracação e $B$ o pátio, ao longo dessa curva $A$ tende a concentrar berço e $B$ a concentrar armazenagem --- o casamento eficiente entre o perfil de cada terminal e o insumo que ele usa com mais intensidade.
\rubricalabel
\begin{itemize}
\item Escreve corretamente $\text{TMS}^A = 2\,x_2^A/x_1^A$ e $\text{TMS}^B = \tfrac12\,x_2^B/x_1^B$: 1{,}5 pt (0{,}75 cada).
\item Impõe a viabilidade $x_1^B = 12 - x_1^A$, $x_2^B = 12 - x_2^A$: 0{,}5 pt.
\item Iguala as $\text{TMS}$ e isola corretamente $x_2^A = \dfrac{4 x_1^A}{16 - x_1^A}$: 2 pt.
\item Verifica a passagem pelos dois cantos (origem de $A$ e origem de $B$): 1 pt (0{,}5 cada).
\item Erro algébrico isolado (ex.: cruzamento de termos correto, mas erro de sinal ao reagrupar) com método íntegro: penalidade leve de $-0{,}5$ pt, sem zerar o item.
\end{itemize}
\end{solucao}

\textbf{b)} (5 pontos) Suponha que a dotação inicial seja $\omega^A = (8,4)$ e $\omega^B = (4,8)$. (i) Mostre que essa alocação já está sobre a curva de contrato (logo, já é Pareto-eficiente). (ii) Determine o preço relativo de equilíbrio walrasiano $p_1^{*}$ (com $p_2=1$) e verifique que, partindo dessa dotação, \emph{não há troca} no equilíbrio. Relacione o resultado ao 1.\textsuperscript{o} e ao 2.\textsuperscript{o} Teoremas do Bem-Estar.
\resolucaobox{5.5cm}
\begin{solucao}
\textbf{Passo 1 --- A dotação está na curva de contrato.} Substituindo $x_1^A = 8$ na equação do item (a): $x_2^A = \dfrac{4\cdot 8}{16 - 8} = \dfrac{32}{8} = 4$. Como $\omega^A = (8,4)$ satisfaz $x_2^A = 4$, a dotação \emph{já} está sobre a curva de contrato e é, portanto, Pareto-eficiente.

\textbf{Passo 2 --- $\text{TMS}$ comum na dotação.} $\text{TMS}^A = 2\cdot\dfrac{4}{8} = 1$ e $\text{TMS}^B = \dfrac12\cdot\dfrac{8}{4} = 1$. As taxas coincidem em $1$ --- confirmação independente da eficiência.

\textbf{Passo 3 --- Preço de equilíbrio e ausência de troca.} O preço que sustenta uma alocação eficiente iguala a $\text{TMS}$ comum: $p_1^{*}/p_2 = \text{TMS} = 1$, logo $p_1^{*} = 1$. Verificação pelas demandas Cobb-Douglas: a renda de $A$ é $m^A = p_1^{*}\,\omega_1^A + \omega_2^A = 1\cdot 8 + 4 = 12$, e
\[ x_1^A = \frac{2}{3}\frac{m^A}{p_1^{*}} = \frac{2}{3}\cdot 12 = 8 = \omega_1^A, \qquad x_2^A = \frac{1}{3}\,m^A = 4 = \omega_2^A. \]
O operador $A$ demanda exatamente sua dotação; por simetria, $B$ também ($m^B=12$, $x_1^B=4$, $x_2^B=8$). Mercados fecham sem troca: a dotação \emph{já é} o equilíbrio walrasiano.

\textbf{Passo 4 --- Conexão com os teoremas.} Pelo 1.\textsuperscript{o} Teorema, o equilíbrio é eficiente --- e de fato a dotação está na curva de contrato. Pelo 2.\textsuperscript{o} Teorema, qualquer alocação eficiente é sustentável por preços com a dotação adequada; aqui a própria dotação coincide com a alocação eficiente, e o preço sustentador é $p_1^{*}=1$, sem necessidade de transferências.

\textbf{Intuição econômica:} quando os terminais já partem de uma divisão em que avaliam os insumos do porto à mesma taxa marginal, não há negócio a fazer --- todo ganho de troca já foi exaurido. O sistema de preços competitivo apenas \emph{ratifica} a divisão existente: ele aponta $p_1^{*}=1$ e cada operador, otimizando, permanece onde está.
\rubricalabel
\begin{itemize}
\item Mostra que $(8,4)$ satisfaz a curva de contrato (substituição correta dando $x_2^A=4$): 1{,}5 pt.
\item Calcula a $\text{TMS}$ comum $=1$ na dotação \emph{ou} argumenta a eficiência de forma equivalente: 1 pt.
\item Obtém $p_1^{*} = 1$ (via $\text{TMS}$ comum ou via clearing das demandas): 1{,}5 pt.
\item Verifica explicitamente a ausência de troca (demanda $=$ dotação) e relaciona ao 1.\textsuperscript{o}/2.\textsuperscript{o} Teoremas: 1 pt.
\item Erro aritmético com método correto (ex.: renda bem montada, conta final trocada): penalidade leve de $-0{,}5$ pt.
\item Concluir $p_1^{*}=1$ sem nenhuma verificação (nem $\text{TMS}$, nem clearing): teto de 3 pts no item.
\end{itemize}
\end{solucao}

\textbf{c)} (5 pontos) Mantenha a economia dos itens anteriores, mas agora a dotação inicial passa a ser $\hat\omega^A = (3,6)$ e $\hat\omega^B = (9,6)$. (i) Verifique que essa dotação \emph{não} está sobre a curva de contrato, resolva o equilíbrio walrasiano (encontre $p_1^{*}$ com $p_2=1$, a alocação final e os fluxos líquidos de troca de cada operador) e \emph{confirme} a alocação via curva de contrato. (ii) Argumente, de forma geral, por que \emph{qualquer} dotação situada sobre a reta orçamentária $\{\,\omega : p^{*}\cdot\omega = p^{*}\cdot\hat\omega^A\,\}$ sustenta exatamente a \emph{mesma} alocação de equilíbrio, e conecte essa invariância ao 2.\textsuperscript{o} Teorema do Bem-Estar.
\resolucaobox{5.5cm}
\begin{solucao}
\textbf{Passo 1 --- A dotação está fora da curva de contrato.} Na curva, $x_1^A=3 \Rightarrow x_2^A = \dfrac{4\cdot 3}{16-3} = \dfrac{12}{13}\approx 0{,}92$. Mas $\hat\omega^A=(3,6)$ tem $x_2^A=6\neq \tfrac{12}{13}$: a dotação está \emph{fora} da curva, logo é ineficiente --- há ganhos de troca a explorar. (De fato $\text{TMS}^A = 2\cdot 6/3 = 4$ enquanto $\text{TMS}^B = \tfrac12\cdot 9/6 = 3/4$; avaliações marginais bem distintas.)

\textbf{Passo 2 --- Equilíbrio walrasiano.} Rendas: $m^A = 3p_1 + 6$, $m^B = 9p_1 + 6$. Demandas pelo bem $1$:
\[ x_1^A = \frac{2}{3}\frac{m^A}{p_1} = \frac{2}{3}\Big(3 + \frac{6}{p_1}\Big) = 2 + \frac{4}{p_1}, \qquad x_1^B = \frac{1}{3}\frac{m^B}{p_1} = \frac{1}{3}\Big(9 + \frac{6}{p_1}\Big) = 3 + \frac{2}{p_1}. \]
Clearing do mercado $1$: $x_1^A + x_1^B = 12 \Rightarrow 5 + \dfrac{6}{p_1} = 12 \Rightarrow \dfrac{6}{p_1} = 7 \Rightarrow p_1^{*} = \dfrac{6}{7}.$

\textbf{Passo 3 --- Alocação final, confirmação e fluxos.} Com $p_1^{*}=\tfrac{6}{7}$: $m^A = 3\cdot\tfrac{6}{7} + 6 = \tfrac{18}{7} + \tfrac{42}{7} = \tfrac{60}{7}$, logo
\[ x_1^A = \frac{2}{3}\frac{m^A}{p_1^{*}} = \frac{2}{3}\cdot\frac{60/7}{6/7} = \frac{2}{3}\cdot 10 = \frac{20}{3}\approx 6{,}67, \qquad x_2^A = \frac{1}{3}\,m^A = \frac{20}{7}\approx 2{,}86. \]
Por viabilidade, $x^B = \big(12-\tfrac{20}{3},\,12-\tfrac{20}{7}\big) = \big(\tfrac{16}{3},\,\tfrac{64}{7}\big)$. \emph{Confirmação pela curva de contrato:} em $x_1^A=\tfrac{20}{3}$, $x_2^A=\dfrac{4\cdot 20/3}{16-20/3}=\dfrac{80/3}{28/3}=\dfrac{80}{28}=\dfrac{20}{7}$ --- bate exatamente, e o 1.\textsuperscript{o} Teorema é satisfeito (equilíbrio na curva de contrato). \emph{Fluxos líquidos de} $A$: de $(3,6)$ para $\big(\tfrac{20}{3},\tfrac{20}{7}\big)$, \emph{compra} $\tfrac{20}{3}-3=\tfrac{11}{3}\approx 3{,}67$ horas-berço e \emph{vende} $6-\tfrac{20}{7}=\tfrac{22}{7}\approx 3{,}14$ lotes de pátio; $B$ faz o espelho. Valor líquido da troca de $A$: $p_1^{*}\cdot\tfrac{11}{3} - p_2\cdot\tfrac{22}{7} = \tfrac{6}{7}\cdot\tfrac{11}{3} - \tfrac{22}{7} = \tfrac{22}{7}-\tfrac{22}{7}=0$, consistente com a restrição orçamentária.

\textbf{Passo 4 --- Invariância sobre a reta orçamentária.} A demanda walrasiana de cada operador depende da dotação \emph{apenas} através de sua renda $m^i = p^{*}\cdot\hat\omega^i$. Em $p^{*}=\big(\tfrac{6}{7},1\big)$, toda dotação de $A$ sobre a reta $p^{*}\cdot\omega = p^{*}\cdot(3,6) = \tfrac{18}{7}+\tfrac{42}{7}=\tfrac{60}{7}$ gera a \emph{mesma} renda $m^A=\tfrac{60}{7}$, logo a \emph{mesma} cesta demandada $\big(\tfrac{20}{3},\tfrac{20}{7}\big)$; e como o total é fixo, $B$ recebe $\big(\tfrac{16}{3},\tfrac{64}{7}\big)$. \emph{Conexão com o 2.\textsuperscript{o} Teorema:} a alocação eficiente $\big(\tfrac{20}{3},\tfrac{20}{7}\big)/\big(\tfrac{16}{3},\tfrac{64}{7}\big)$ pode ser descentralizada como equilíbrio walrasiano a partir de \emph{qualquer} ponto dessa reta de preços --- o que a autoridade portuária controla, via redistribuição de dotações de mesmo valor ao preço $p^{*}$, é apenas \emph{quem paga a quem} (o tamanho dos fluxos), não a alocação final eficiente, que permanece a mesma.

\textbf{Intuição econômica:} preços competitivos só enxergam o \emph{valor} da dotação, não sua composição. Dois pontos de partida com o mesmo valor a preços de equilíbrio são economicamente indistinguíveis: o mercado conduz ambos ao mesmo destino eficiente, ajustando apenas o tamanho das trocas. É a tradução operacional do 2.\textsuperscript{o} Teorema --- separar a questão distributiva (onde colocar a dotação sobre a reta) da questão de eficiência (o destino), desde que as preferências sejam convexas, como aqui.
\rubricalabel
\begin{itemize}
\item Verifica que $\hat\omega^A=(3,6)$ está fora da curva de contrato (ou $\text{TMS}^A\neq\text{TMS}^B$ na dotação): 0{,}5 pt.
\item Monta as demandas Cobb-Douglas com rendas corretas e resolve o clearing para $p_1^{*}=\tfrac{6}{7}$: 1{,}5 pt.
\item Obtém a alocação final $\big(\tfrac{20}{3},\tfrac{20}{7}\big)/\big(\tfrac{16}{3},\tfrac{64}{7}\big)$, confirma pela curva de contrato e dá os fluxos líquidos ($+\tfrac{11}{3}$ berço, $-\tfrac{22}{7}$ pátio para $A$): 1 pt.
\item Argumenta a invariância: demanda depende da dotação só via renda $p^{*}\cdot\omega$, logo toda dotação sobre a reta orçamentária dá a mesma alocação: 1{,}5 pt.
\item Conecta explicitamente ao 2.\textsuperscript{o} Teorema (descentralização / separação eficiência $\times$ distribuição): 0{,}5 pt.
\item Erro aritmético com método correto (ex.: demandas certas, conta final do clearing trocada): penalidade leve de $-0{,}5$ pt.
\item Resolver o equilíbrio corretamente mas sem o argumento geral de invariância (ii): teto de 3 pts no item.
\end{itemize}
\end{solucao}


\blocotitulo{Bloco 3 --- Externalidades e Bens Públicos \hfill \normalsize (25 pontos)}

\textbf{Questão 5.} (8 pontos) \textbf{Itens objetivos.} Nos itens \textbf{a)} e \textbf{b)}, assinale a \emph{única} alternativa correta (3 pontos cada; não há penalização por erro). No item \textbf{c)}, responda \textbf{Verdadeiro ou Falso} (2 pontos; vale a regra de anulação descrita na capa). Registre suas respostas no quadro-resposta ao final do bloco.

\textbf{a)} (3 pontos) A polinização por abelhas é um caso clássico de externalidade positiva: na região de Limeira (SP), polo citrícola, a instalação de colmeias por apicultores beneficia os laranjais vizinhos. Suponha que o benefício marginal privado de instalar colmeias seja dado pela demanda inversa $P(q)=30-3q$ (em R\$~mil por colmeia), que o custo marginal privado do apicultor seja constante, $\text{MPC}=9$, e que cada colmeia gere um benefício externo marginal constante de R\$~6 mil aos laranjais. Qual é o subsídio pigouviano ótimo por colmeia que conduz o apicultor à instalação socialmente eficiente?

\begin{enumerate}[label=\Alph*., leftmargin=3.2em, itemsep=0.05em, topsep=0.1em]
  \item $9$
  \item $6$
  \item $2$
  \item $3$
  \item $12$
\end{enumerate}
\gabaritoitem{B}
\begin{solucao}
\textbf{Passo 1 — o que se pede.} Numa externalidade \emph{positiva}, o benefício social de cada colmeia excede o privado pelo benefício externo marginal ($\text{MEB}$) que recai sobre os laranjais. O mercado privado provê pouco; o subsídio pigouviano ótimo fecha exatamente a diferença, igualando $\text{MEB}$ na margem eficiente.

\textbf{Passo 2 — equilíbrio privado.} Deixado sozinho, o apicultor instala colmeias até que seu benefício marginal privado iguale seu custo marginal privado:
\begin{align*}
P(q)=\text{MPC}\;\Longrightarrow\; 30-3q=9\;\Longrightarrow\; q^{\text{priv}}=7.
\end{align*}

\textbf{Passo 3 — ótimo social e subsídio.} O ótimo social iguala o \emph{benefício marginal social} ($P(q)+\text{MEB}$) ao custo marginal:
\begin{align*}
P(q)+\text{MEB}=\text{MPC}\;\Longrightarrow\; (30-3q)+6=9\;\Longrightarrow\; q^{*}=9.
\end{align*}
O subsídio pigouviano ótimo é o benefício externo marginal avaliado no ótimo. Como $\text{MEB}=6$ é constante, $s^{*}=6$. Verificação: com o subsídio, o apicultor maximiza $P(q)+s^{*}=\text{MPC}$, isto é $30-3q=9-6=3\Rightarrow q=9=q^{*}$. O subsídio empurra a instalação de $7$ para $9$ colmeias.

\textbf{Passo 4 — distratores tentadores.} A alternativa $9$ é a mais sedutora por dupla razão: é o \emph{número de colmeias} no ótimo social ($q^{*}=9$, confusão da quantidade-alvo com o instrumento) e coincide com o $\text{MPC}=9$, valor saliente no enunciado — dois erros desembocam no mesmo número. A alternativa $2$ corresponde à \emph{variação} de quantidade $q^{*}-q^{\text{priv}}=9-7=2$, erro de identificar o subsídio com o deslocamento de output em vez do diferencial de benefício marginal. A alternativa $3$ é o preço $P(q^{*})=30-3\cdot 9=3$ no ótimo (o que o apicultor paga líquido após o subsídio), não o subsídio. A alternativa $12=2\,\text{MEB}$ é um erro de escala (dobrar a externalidade).

\textbf{Intuição econômica:} cada colmeia gera um ganho de R\$~6 mil que o apicultor não captura — ele poliniza os laranjais do vizinho de graça. Por isso instala colmeias de menos: ignora R\$~6 mil de benefício social por colmeia. O subsídio de R\$~6 mil por colmeia internaliza esse ganho externo, alinhando o incentivo privado ao social e elevando a instalação ao nível eficiente. O valor do subsídio é o tamanho da externalidade na margem, não a quantidade nem sua variação.
\end{solucao}

\textbf{b)} (3 pontos) A CETESB regula a poluição industrial no polo de Cubatão (SP). Suponha que, para certa indústria, o benefício marginal de emitir seja $\text{MB}(e)=20-2e$ e que o dano marginal imposto a uma colônia de pescadores a jusante seja $\text{MD}(e)=2e$, com $e$ em toneladas/dia e valores em R\$~mil. Os direitos sobre o rio cabem aos pescadores (direito a água limpa) e os custos de transação são nulos. Pela barganha de Coase, qual é a compensação mínima que a indústria precisa pagar aos pescadores para emitir o nível eficiente?

\begin{enumerate}[label=\Alph*., leftmargin=3.2em, itemsep=0.05em, topsep=0.1em]
  \item $5$
  \item $50$
  \item $75$
  \item $25$
  \item $100$
\end{enumerate}
\gabaritoitem{D}
\begin{solucao}
\textbf{Passo 1 — o que se pede.} Com direitos atribuídos aos pescadores, o ponto de partida é emissão zero (eles têm direito a água limpa). A indústria só polui se pagar uma compensação aceitável. A barganha de Coase leva ao nível eficiente; pede-se o \emph{piso} dessa compensação — quanto a indústria precisa, no mínimo, pagar para que os pescadores aceitem sair de $e=0$ até o eficiente.

\textbf{Passo 2 — nível eficiente.} O ótimo iguala benefício e dano marginais:
\begin{align*}
\text{MB}(e)=\text{MD}(e)\;\Longrightarrow\; 20-2e=2e\;\Longrightarrow\; e^{*}=5\ \text{ton/dia}.
\end{align*}

\textbf{Passo 3 — piso e teto da barganha (curvas distintas).} O \emph{piso} é a compensação mínima que os pescadores exigem para aceitar a poluição: igual ao \emph{dano total} que sofrem ao passar de $0$ a $e^{*}=5$, que é a \emph{área sob $\text{MD}$}:
\begin{align*}
\text{piso}=\int_{0}^{5}\text{MD}(e)\,\mathrm{d}e=\int_{0}^{5} 2e\,\mathrm{d}e=\big[e^{2}\big]_{0}^{5}=25.
\end{align*}
O \emph{teto} é a máxima disposição a pagar da indústria: igual ao \emph{benefício total} de poluir até $5$, que é a \emph{área sob $\text{MB}$} (curva diferente):
\begin{align*}
\text{teto}=\int_{0}^{5}\text{MB}(e)\,\mathrm{d}e=\int_{0}^{5}(20-2e)\,\mathrm{d}e=\big[20e-e^{2}\big]_{0}^{5}=75.
\end{align*}
Note que cada lado integra uma curva própria: pescadores olham $\text{MD}$ (o que perdem), indústria olha $\text{MB}$ (o que ganha). Qualquer pagamento em $[25,75]$ viabiliza a troca; o \emph{mínimo} é R\$~25 mil/dia. O excedente social criado, $75-25=50$, é o que as partes repartem conforme o poder de barganha.

\textbf{Passo 4 — distratores tentadores.} A alternativa $75$ é o \emph{teto} — o benefício total da indústria (área sob $\text{MB}$), máximo que ela aceitaria pagar, confundido com o mínimo que precisa pagar. A alternativa $50$ é o \emph{excedente social líquido} ($75-25$), o ganho total da barganha, não a transferência aos pescadores. A alternativa $100=\int_{0}^{10}2e\,\mathrm{d}e$ é o dano (área sob $\text{MD}$) avaliado até a emissão não regulada $e=10$ (onde $\text{MB}=0$), erro de integrar $\text{MD}$ até a quantidade errada. A alternativa $5$ é o próprio $e^{*}$, quantidade e não valor.

\textbf{Intuição econômica:} o Teorema de Coase garante que, com direitos definidos e sem custos de transação, as partes negociam até o nível eficiente $e^{*}=5$, independentemente de quem detém o direito. O que muda com a alocação de direitos é a \emph{distribuição}: como o direito é dos pescadores, a indústria compra o direito de poluir e os compensa por seu prejuízo. O piso da compensação é o dano que eles efetivamente sofrem — a área sob $\text{MD}$, R\$~25 mil/dia. Abaixo disso, recusam; acima de R\$~75 mil (área sob $\text{MB}$), a indústria desiste. Entre os dois, o acordo é mutuamente vantajoso.
\end{solucao}

\textbf{c)} (2 pontos, V ou F) A iluminação pública municipal, custeada pela taxa de iluminação cobrada na conta de energia, é um bem público local: todos os moradores de uma área desfrutam simultaneamente do mesmo nível de iluminação $G$. \textbf{Afirmação:} num equilíbrio de Lindahl que financie o nível eficiente de iluminação, o que a condição de equilíbrio exige é que \emph{cada morador pague o mesmo preço de Lindahl por unidade de $G$}, pois só assim a arrecadação cobre o custo marginal da iluminação.

\gabaritoitem{F}
\begin{solucao}
\textbf{Passo 1 — o que se afirma.} A sentença sustenta que o equilíbrio de Lindahl \emph{exige preços iguais} e que seria essa uniformidade que faz a arrecadação cobrir o custo. Ambas as partes estão erradas: os preços de Lindahl são \emph{personalizados}, e o que cobre o custo marginal é a \emph{soma} desses preços — não a sua igualdade.

\textbf{Passo 2 — o que realmente caracteriza Lindahl.} Num equilíbrio de Lindahl, cada consumidor $i$ enfrenta um preço-cota próprio $\tau^{i}$ por unidade de $G$ e, tomando-o como dado, demanda voluntariamente uma quantidade do bem público. O equilíbrio exige que (i) todos demandem o \emph{mesmo} nível $G^{*}$ — a quantidade é comum porque o bem é não-rival — e (ii) os preços personalizados \emph{somem} ao custo marginal: $\sum_i \tau^{i}=\text{TMT}_{G,x}$. A condição individual de ótimo iguala o preço de cada um à sua \emph{taxa marginal de substituição}: $\tau^{i}=\text{TMS}^{i}_{G,x}(G^{*})$. É a \emph{soma} $\sum_i\tau^i$ — e não um preço comum — que custeia a unidade marginal de $G$.

\textbf{Passo 3 — por que a afirmação falha.} Se os moradores têm valorações distintas pela iluminação, suas $\text{TMS}^{i}$ em $G^{*}$ diferem; logo seus preços de Lindahl diferem. Concretamente, com $u^{i}=x^{i}+\theta_i\ln G$ e $\text{TMT}=1$, Samuelson dá $G^{*}=\sum_i\theta_i$ e $\tau^{i}=\theta_i/G^{*}$: a soma $\sum_i\tau^i=\sum_i\theta_i/G^{*}=1=\text{TMT}$ cobre o custo, mas os termos individuais são desiguais — quem valoriza mais paga mais por unidade. Impor preço uniforme $\tau^i=1/n$ quebraria a igualdade $\tau^i=\text{TMS}^i$ para todos exceto no caso degenerado de moradores idênticos. Como nada no enunciado impõe homogeneidade, a afirmação é \textbf{falsa}.

\textbf{Intuição econômica:} a não-rivalidade impõe que todos ``consumam'' o mesmo nível de iluminação — a luz acesa ilumina a rua inteira de uma vez. O que cobre o custo de mais uma unidade não é todo mundo pagar o \emph{mesmo} preço, e sim a \emph{soma} de preços personalizados igualar o custo marginal. A eficiência exige justamente preços \emph{diferentes}: quem se beneficia mais paga uma fração maior. Confundir ``soma dos preços cobre o custo'' com ``cada um paga o mesmo preço'' é trocar a condição de Lindahl ($\sum_i\tau^i=\text{TMT}$) por uma cobrança per capita uniforme, que não implementa o ótimo.
\end{solucao}

\ifsolucao\else
\par\vspace{0.5em}
\begin{center}
\textbf{Quadro-resposta do Bloco 3 --- escreva a alternativa (A--E) ou V/F de cada item:}\\[0.6em]
\renewcommand{\arraystretch}{1.45}
\begin{tabular}{|>{\centering\arraybackslash}p{2.6cm}|>{\centering\arraybackslash}p{1.4cm}|>{\centering\arraybackslash}p{1.4cm}|>{\centering\arraybackslash}p{1.4cm}|}
  \hline
  \textbf{Item} & a) & b) & c) \\ \hline
  \textbf{Resposta} &  &  &  \\ \hline
\end{tabular}
\end{center}
{\small\itshape Atenção: apenas o quadro acima será considerado na correção dos itens objetivos deste bloco.}\par\vspace{0.6em}
\fi

\textbf{Questão 6.} (17 pontos) A Votorantim Cimentos opera uma fábrica de cimento no interior de São Paulo. A produção de clínquer libera material particulado que se deposita sobre a área urbana vizinha, gerando custos de saúde e limpeza que recaem sobre os moradores — uma externalidade negativa de produção que a firma não interioriza. Modele a fábrica como uma firma que escolhe o nível de produção $q$ (em milhares de toneladas/mês) para maximizar seu lucro $\pi(q)=R(q)-C(q)$, com receita $R(q)=42q-q^{2}$ e custo de produção $C(q)=6q$ (valores em R\$~milhões/mês). O dano externo imposto à população é $D(q)=q^{2}/2$, também em R\$~milhões/mês, de modo que o dano marginal é $\text{MD}(q)=q$. Trate $q$ como contínuo e estritamente positivo, e suponha que o regulador (CETESB) busque maximizar o excedente social $\pi(q)-D(q)$.

\textbf{a)} (6 pontos) Determine, por condição de primeira ordem, o nível de produção privado $q^{\text{priv}}$ (que a firma escolhe ignorando o dano) e o nível socialmente eficiente $q^{*}$ (que maximiza $\pi(q)-D(q)$). Explique economicamente por que $q^{\text{priv}}>q^{*}$ e o que separa as duas decisões na margem.
\resolucaobox{5.5cm}
\begin{solucao}
\textbf{Passo 1 — benefício privado marginal.} O lucro privado é $\pi(q)=R(q)-C(q)=(42q-q^{2})-6q=36q-q^{2}$. A firma, ignorando o dano externo, maximiza $\pi$:
\begin{align*}
\pi'(q)=R'(q)-C'(q)=(42-2q)-6=36-2q.
\end{align*}

\textbf{Passo 2 — produção privada.} A CPO privada zera o lucro marginal:
\begin{align*}
36-2q=0\;\Longrightarrow\; q^{\text{priv}}=18.
\end{align*}
A segunda ordem confirma máximo: $\pi''=-2<0$.

\textbf{Passo 3 — produção eficiente.} O regulador maximiza o excedente social $W(q)=\pi(q)-D(q)=36q-q^{2}-\tfrac{1}{2}q^{2}$. A CPO social subtrai o dano marginal $\text{MD}(q)=D'(q)=q$ do lucro marginal:
\begin{align*}
W'(q)=\underbrace{(36-2q)}_{\pi'(q)}-\underbrace{q}_{\text{MD}(q)}=36-3q=0\;\Longrightarrow\; q^{*}=12.
\end{align*}
Segunda ordem: $W''=-3<0$, máximo. Logo $q^{\text{priv}}=18>q^{*}=12$.

\textbf{Passo 4 — por que $q^{\text{priv}}>q^{*}$.} As duas decisões compartilham o lucro marginal $\pi'(q)=36-2q$; o que as separa é exatamente o dano marginal $\text{MD}(q)=q$, que entra na conta social mas não na privada. A firma produz até o ponto em que seu lucro marginal se anula ($q=18$); o regulador para antes, em $q=12$, porque a partir daí o lucro marginal de mais uma unidade já não cobre o dano que ela impõe aos moradores. Entre $12$ e $18$ unidades, a firma ainda lucra na margem ($\pi'>0$), mas cada unidade extra custa à sociedade mais do que rende — o lucro marginal é menor que o dano marginal.

\textbf{Intuição econômica:} a poluição é um custo real que não aparece na contabilidade da fábrica. Como a Votorantim Cimentos não paga pelo material particulado que deposita sobre a cidade, ela produz como se esse custo não existisse e empurra a produção além do que é socialmente desejável. O excesso $q^{\text{priv}}-q^{*}=6$ mil toneladas/mês é o sobre-output característico de toda externalidade negativa de produção: o mercado descentralizado iguala custo \emph{privado} ao benefício marginal, não o custo \emph{social}.
\rubricalabel
\begin{itemize}
\item Escreve $\pi(q)=36q-q^{2}$ e obtém o lucro marginal $\pi'(q)=36-2q$: 1{,}5 pt.
\item Resolve a CPO privada e obtém $q^{\text{priv}}=18$: 1{,}5 pt.
\item Monta a CPO social subtraindo $\text{MD}(q)=q$ do lucro marginal e obtém $q^{*}=12$: 2 pts.
\item Justifica economicamente $q^{\text{priv}}>q^{*}$ pelo dano marginal não-interiorizado (lucro marginal positivo, mas menor que o dano, entre $12$ e $18$): 1 pt.
\item Erro aritmético isolado com método integralmente correto (ex.: monta $36-3q=0$ e erra a divisão): penalidade leve de $-0{,}5$ pt (não zera o item).
\item Esquecer de subtrair o dano marginal na CPO social (resolver as duas como se fossem o problema privado): teto de 2 pts no item.
\end{itemize}
\end{solucao}

\textbf{b)} (6 pontos) Para corrigir a externalidade, a CETESB cobra um imposto pigouviano $t$ por unidade produzida, de modo que a firma passa a maximizar $\pi(q)-tq$. Determine a alíquota ótima $t^{*}$ e \emph{verifique} que, enfrentando esse imposto, a firma escolhe voluntariamente exatamente $q^{*}$. Calcule também a arrecadação do imposto e o dano total no ótimo, e comente se eles coincidem.
\resolucaobox{5.5cm}
\begin{solucao}
\textbf{Passo 1 — princípio do imposto pigouviano.} O imposto pigouviano ótimo iguala o dano marginal avaliado no nível eficiente, $t^{*}=\text{MD}(q^{*})$. Como $\text{MD}(q)=q$ e $q^{*}=12$:
\begin{align*}
t^{*}=\text{MD}(q^{*})=q^{*}=12\ \text{(R\$~milhões por mil toneladas)}.
\end{align*}

\textbf{Passo 2 — problema da firma com imposto.} Com o imposto, a firma maximiza $\pi(q)-t^{*}q=36q-q^{2}-12q=24q-q^{2}$. A CPO é
\begin{align*}
\frac{\mathrm{d}}{\mathrm{d}q}\big(\pi(q)-t^{*}q\big)=\pi'(q)-t^{*}=(36-2q)-12=0.
\end{align*}

\textbf{Passo 3 — verificação.} Resolvendo:
\begin{align*}
36-2q=12\;\Longrightarrow\; 2q=24\;\Longrightarrow\; q=12=q^{*}.\ \checkmark
\end{align*}
A firma, agindo no próprio interesse, escolhe exatamente o nível eficiente: o imposto faz o lucro marginal privado bater com o dano marginal justamente em $q^{*}=12$. Segunda ordem inalterada ($-2<0$), máximo.

\textbf{Passo 4 — arrecadação versus dano total.} A arrecadação no ótimo é $t^{*}q^{*}=12\cdot 12=144$. O dano total é $D(q^{*})=\tfrac{1}{2}(q^{*})^{2}=\tfrac{1}{2}\cdot 144=72$. Eles \emph{não} coincidem: a arrecadação ($144$) excede o dano total ($72$). A razão é que o imposto pigouviano precisa igualar o dano \emph{marginal} ($\text{MD}(q^{*})=12$), não o dano \emph{médio}. Com dano marginal crescente, $\text{MD}(q^{*})=12$ está acima do dano médio $D(q^{*})/q^{*}=6$; multiplicado por todas as unidades, gera receita maior que o dano agregado. O objetivo do imposto é \emph{corrigir a margem}, não ``pagar exatamente a conta'' do dano total.

\textbf{Intuição econômica:} o imposto pigouviano funciona como um preço que a firma é obrigada a pagar pelo dano que cada tonelada extra causa. Ao confrontar a Votorantim Cimentos com o custo marginal que ela impunha de graça à cidade, o regulador faz a firma internalizar a externalidade sem precisar ditar a quantidade: basta acertar o preço da poluição na margem. O fato de a arrecadação superar o dano total é normal quando o dano cresce com a escala — e ilustra que ``imposto pigouviano'' não é sinônimo de ``ressarcimento integral do prejuízo'', mas de correção do incentivo marginal.
\rubricalabel
\begin{itemize}
\item Identifica o princípio $t^{*}=\text{MD}(q^{*})$ e calcula $t^{*}=12$: 1{,}5 pt.
\item Monta o problema da firma $\pi(q)-t^{*}q$ e a CPO $\pi'(q)=t^{*}$: 1{,}5 pt.
\item Verifica que a firma escolhe $q=12=q^{*}$ sob o imposto: 1{,}5 pt.
\item Calcula arrecadação $t^{*}q^{*}=144$ e dano total $D(q^{*})=72$ e comenta que não coincidem (margem vs.\ média): 1{,}5 pt.
\item Erro aritmético isolado com método correto: penalidade leve de $-0{,}5$ pt.
\item Fixar $t^{*}$ pelo dano marginal na produção \emph{privada} ($\text{MD}(18)=18$) em vez de na eficiente: teto de 3 pts no item.
\end{itemize}
\end{solucao}

\textbf{c)} (5 pontos) Suponha que o regulador nada faça e a firma produza $q^{\text{priv}}$. Calcule a perda de peso-morto (\emph{deadweight loss}) dessa inação — isto é, o excedente social perdido por produzir $q^{\text{priv}}$ em vez de $q^{*}$. Apresente o cálculo como a área entre o lucro marginal e o dano marginal no intervalo relevante e interprete o resultado.
\resolucaobox{5.5cm}
\begin{solucao}
\textbf{Passo 1 — o que é a perda de peso-morto aqui.} Para cada unidade produzida \emph{além} de $q^{*}=12$, o dano marginal supera o lucro marginal: a unidade rende menos à firma do que custa à sociedade. A perda social acumulada de ir de $q^{*}=12$ até $q^{\text{priv}}=18$ é a integral dessa diferença:
\begin{align*}
\text{DWL}=\int_{q^{*}}^{q^{\text{priv}}}\big[\text{MD}(q)-\pi'(q)\big]\,\mathrm{d}q=\int_{12}^{18}\big[q-(36-2q)\big]\,\mathrm{d}q=\int_{12}^{18}(3q-36)\,\mathrm{d}q.
\end{align*}

\textbf{Passo 2 — cálculo da integral.}
\begin{align*}
\int_{12}^{18}(3q-36)\,\mathrm{d}q=\Big[\tfrac{3}{2}q^{2}-36q\Big]_{12}^{18}=\big(486-648\big)-\big(216-432\big)=(-162)-(-216)=54.
\end{align*}
A perda de peso-morto é R\$~54 milhões/mês.

\textbf{Passo 3 — verificação pela área do triângulo.} O integrando $3q-36$ é o \emph{negativo} do excedente social marginal $W'(q)=36-3q$, que vale zero em $q^{*}=12$ e $-18$ em $q^{\text{priv}}=18$. A região é um triângulo de base $q^{\text{priv}}-q^{*}=6$ e altura $|W'(18)|=18$:
\begin{align*}
\text{DWL}=\tfrac{1}{2}\cdot 6\cdot 18=54.\ \checkmark
\end{align*}
Os dois métodos batem.

\textbf{Passo 4 — interpretação.} As $6$ mil toneladas/mês produzidas em excesso geram, somadas, R\$~54 milhões de perda líquida: é o tamanho do prejuízo social que o imposto pigouviano do item anterior eliminaria ao recuar a produção de $18$ para $12$. Note que a DWL \emph{não} é o dano total ($72$) nem a arrecadação ($144$): é apenas a parcela de excedente \emph{perdida} pela escolha errada de quantidade — o triângulo entre as curvas de lucro marginal e dano marginal sobre o sobre-output.

\textbf{Intuição econômica:} deixar a fábrica produzir livremente não custa à sociedade todo o dano da poluição — afinal, parte dessa produção vale a pena mesmo com o dano. O que se perde é só o ``excesso'': as toneladas finais, cujo lucro já não compensa o estrago que causam. Esse triângulo de R\$~54 milhões é exatamente o ganho de bem-estar que a CETESB captura ao tributar a margem; é a medida monetária da ineficiência da inação, e a razão econômica para regular.
\rubricalabel
\begin{itemize}
\item Identifica a DWL como a integral de $\text{MD}(q)-\pi'(q)$ entre $q^{*}=12$ e $q^{\text{priv}}=18$ (área entre as curvas sobre o sobre-output): 1{,}5 pt.
\item Monta o integrando $3q-36$ corretamente: 1 pt.
\item Calcula a integral e obtém $\text{DWL}=54$ (ou pela área do triângulo $\tfrac12\cdot 6\cdot 18=54$): 1{,}5 pt.
\item Interpreta que a DWL é o excedente perdido pelo sobre-output, distinto do dano total e da arrecadação: 1 pt.
\item Erro aritmético isolado com método correto (limites e integrando certos, conta final trocada): penalidade leve de $-0{,}5$ pt.
\item Calcular ``dano total'' $D(18)$ ou $D(12)$ e chamar de DWL (confundir nível com perda marginal acumulada): teto de 2 pts no item.
\end{itemize}
\end{solucao}


\blocotitulo{Bloco 4 --- Assimetria Informacional e Sinalização \hfill \normalsize (25 pontos)}

\textbf{Questão 7.} (8 pontos) \textbf{Itens objetivos.} Nos itens \textbf{a)} e \textbf{b)}, assinale a \emph{única} alternativa correta (3 pontos cada; não há penalização por erro). No item \textbf{c)}, responda \textbf{Verdadeiro ou Falso} (2 pontos; vale a regra de anulação descrita na capa). Registre suas respostas no quadro-resposta ao final do bloco.

\textbf{a)} (3 pontos) Uma fintech de crédito (estilo Serasa Crédito) concede empréstimos de R\$~1 a empreendedores observacionalmente idênticos, metade de cada um de dois tipos. O projeto cauteloso paga $X_S=1{,}2$ em caso de sucesso (probabilidade $p_S=1$); o arrojado paga $X_H=2{,}4$ (probabilidade $p_R=\tfrac12$). O empréstimo é quitado, ao fator bruto $R$, apenas em caso de sucesso; em fracasso, paga-se $0$. O tomador aceita o crédito apenas se seu lucro esperado $p(X-R)$ for não-negativo. A fintech, com custo bruto de captação $1$, fixa um único fator $R$ para todos os tomadores (taxa uniforme) e exige lucro esperado nulo, dado quem efetivamente aceita o crédito. Qual o fator bruto de equilíbrio $R^{*}$?

\begin{enumerate}[label=\Alph*., leftmargin=3.2em, itemsep=0.05em, topsep=0.1em]
  \item $R^{*}=1$
  \item $R^{*}=1{,}2$
  \item $R^{*}=\dfrac{4}{3}$
  \item $R^{*}=2{,}4$
  \item $R^{*}=2$
\end{enumerate}
\gabaritoitem{E}
\begin{solucao}
\textbf{Passo 1 --- o que se pede.} Trata-se de sele\c{c}\~ao adversa em cr\'edito com a taxa de juros como \emph{seletor}: \`a medida que $R$ sobe, os tomadores mais seguros saem do mercado, piorando a composi\c{c}\~ao do conjunto que toma cr\'edito. Pede-se o $R^{*}$ que zera o lucro esperado da fintech, dado quem efetivamente toma o cr\'edito.

\textbf{Passo 2 --- quem aceita o cr\'edito.} O tomador aceita se $p(X-R)\ge 0$, ou seja, se $R\le X$. Assim:
\begin{itemize}
\item Cauteloso ($X_S=1{,}2$): aceita enquanto $R\le 1{,}2$.
\item Arrojado ($X_H=2{,}4$): aceita enquanto $R\le 2{,}4$.
\end{itemize}
O cauteloso \'e o primeiro a desistir quando $R$ sobe --- eis a sele\c{c}\~ao adversa.

\textbf{Passo 3 --- testar a taxa uniforme com os dois tipos.} Se ambos tomam o cr\'edito, a probabilidade m\'edia de quita\c{c}\~ao \'e $\bar p=\tfrac{1+\tfrac12}{2}=0{,}75$. O retorno esperado por empr\'estimo \'e $\bar p\,R=0{,}75R$; o lucro nulo exigiria $0{,}75R=1\Rightarrow R=\tfrac{4}{3}\approx 1{,}33$. Mas $1{,}33>1{,}2$: a esse fator o cauteloso \emph{n\~ao} aceita. Logo um equil\'ibrio com os dois tipos no mercado \'e \textbf{inconsistente}.

\textbf{Passo 4 --- equil\'ibrio com apenas os arrojados.} Para $R\in(1{,}2,\,2{,}4]$ s\'o os arrojados tomam cr\'edito, com $p_R=\tfrac12$. Lucro nulo: $\tfrac12 R=1\Rightarrow R^{*}=2$. Verifica\c{c}\~ao: $2\in(1{,}2,\,2{,}4]$, ent\~ao de fato s\'o os arrojados aceitam ($1{,}2-2<0$) e a fintech zera o lucro ($\tfrac12\cdot 2=1$). Equil\'ibrio consistente e \'unico. O fator bruto $R^{*}=2$ corresponde a juros de $100\%$, com os cautelosos racionados para fora.

\textbf{Distratores tentadores.} O distrator $\tfrac{4}{3}$ \'e o erro de quem aplica a m\'edia \emph{ingenuamente}: calcula $0{,}75R=1$ sem checar que a esse $R$ o cauteloso j\'a abandonou o mercado --- ignora a endogeneidade da composi\c{c}\~ao do conjunto de tomadores, que \'e o cora\c{c}\~ao do problema. O distrator $1{,}2$ confunde o \emph{limiar de participa\c{c}\~ao} do cauteloso (o $R$ m\'aximo que ele aceitaria) com o $R$ de equil\'ibrio; \'e o ponto em que o cauteloso sai, n\~ao onde a fintech zera lucro.

\textbf{Intui\c{c}\~ao econ\^omica:} a taxa de juros n\~ao \'e apenas o pre\c{c}o do cr\'edito --- ela tria os tomadores. Subir os juros para cobrir a inadimpl\^encia expulsa justamente os bons pagadores, deixando no mercado os de maior risco, o que for\c{c}a os juros ainda mais para cima. O equil\'ibrio se assenta onde restam s\'o os arrojados, pagando $100\%$, e o cr\'edito barato que os cautelosos mereceriam simplesmente n\~ao existe.
\end{solucao}

\textbf{b)} (3 pontos) A Magazine Luiza oferece garantia estendida em duas vers\~oes: B\'asica (1 ano de cobertura, $q_1=1$) e Premium (2 anos, $q_2=2$). H\'a dois perfis de cliente em iguais propor\c{c}\~oes, que valoram a cobertura por $\theta q$: o cliente leve ($\theta_L=20$) e o cliente intenso ($\theta_H=30$). A loja n\~ao observa o perfil e precifica o menu para que cada um se autosselecione, extraindo o m\'aximo de receita. Ela cobra do leve seu valor cheio pela B\'asica, $T_1=\theta_L q_1$. Para que o cliente intenso prefira a Premium \`a B\'asica (compatibilidade de incentivo \emph{binding}), qual deve ser o pre\c{c}o $T_2$ da Premium?

\begin{enumerate}[label=\Alph*., leftmargin=3.2em, itemsep=0.05em, topsep=0.1em]
  \item $T_2=R\$~50$
  \item $T_2=R\$~60$
  \item $T_2=R\$~40$
  \item $T_2=R\$~70$
  \item $T_2=R\$~45$
\end{enumerate}
\gabaritoitem{A}
\begin{solucao}
\textbf{Passo 1 --- o que se pede.} \'E um problema de \emph{screening} de segundo grau (menu autosseletivo). A loja fixa $(q_1,T_1)$ e $(q_2,T_2)$ de modo que cada tipo escolha o pacote desenhado para ele. Pede-se $T_2$ que torna a restri\c{c}\~ao de incentivo do tipo intenso exatamente \emph{binding}.

\textbf{Passo 2 --- pre\c{c}o da B\'asica.} O leve recebe $T_1=\theta_L q_1=20\cdot 1=20$, ficando com excedente nulo: $\theta_L q_1-T_1=0$. (O tipo baixo \'e espremido at\'e o excedente zero --- padr\~ao em screening.)

\textbf{Passo 3 --- restri\c{c}\~ao de incentivo do intenso.} O intenso s\'o escolhe a Premium se obtiver dela pelo menos o excedente que teria desviando para a B\'asica:
\[
\theta_H q_2-T_2\;\ge\;\theta_H q_1-T_1.
\]
No ponto \emph{binding} (loja extrai o m\'aximo deixando o desvio apenas empatado):
\[
T_2=\theta_H q_2-\big(\theta_H q_1-T_1\big)=30\cdot 2-\big(30\cdot 1-20\big)=60-10=50.
\]

\textbf{Passo 4 --- verifica\c{c}\~ao das duas autosseles.} Intenso: excedente na Premium $=60-50=10$; na B\'asica $=30-20=10$ --- indiferente, ent\~ao fica na Premium (IC binding, satisfeita). O termo $10$ \'e a \emph{renda informacional} que a loja abre m\~ao para evitar o desvio. Leve: excedente na B\'asica $=0$; na Premium $=\theta_L q_2-T_2=40-50=-10<0$ --- jamais migra (satisfeita). O menu separa os tipos.

\textbf{Distratores tentadores.} O distrator $R\$~60$ \'e o erro de quem cobra do intenso seu \emph{valor cheio} pela Premium ($\theta_H q_2=60$), ignorando a renda informacional: a esse pre\c{c}o o intenso desviaria para a B\'asica, onde teria excedente $30-20=10>0$. O distrator $R\$~40$ aplica \`a Premium a valora\c{c}\~ao do tipo \emph{errado} ($\theta_L q_2=40$), confundindo quem \'e o destinat\'ario do pacote superior.

\textbf{Intui\c{c}\~ao econ\^omica:} quando a loja n\~ao enxerga o perfil, n\~ao consegue extrair tudo de todos. Para impedir que o cliente intenso se ``disfarce'' comprando a B\'asica barata, ela tem de lhe deixar uma fatia de excedente --- a renda informacional. O pre\c{c}o da Premium cai de $60$ para $50$ exatamente nessa fatia. \'E o pre\c{c}o que a assimetria de informa\c{c}\~ao cobra de quem vende: o vendedor sacrifica receita do tipo alto para mant\^e-lo no pacote certo.
\end{solucao}

\textbf{c)} (2 pontos, V ou F) Um sistema de aloca\c{c}\~ao de candidatos a vagas universit\'arias usa aceita\c{c}\~ao diferida (algoritmo de Gale--Shapley), respeitando as prefer\^encias dos candidatos e as prioridades dos cursos. Considere dois candidatos, Ana e Bruno, e duas vagas, Medicina e Direito (uma cada). As prefer\^encias dos candidatos s\~ao: Ana e Bruno \emph{ambos} preferem Medicina a Direito. As prioridades dos cursos s\~ao: Medicina prioriza Ana sobre Bruno; Direito prioriza Bruno sobre Ana. Foi proposta, fora do algoritmo, a aloca\c{c}\~ao $\mu$: Ana $\to$ Direito e Bruno $\to$ Medicina. Afirma-se: \emph{o par (Ana, Medicina) \'e um par bloqueante, de modo que $\mu$ n\~ao \'e est\'avel}.

\gabaritoitem{V}
\begin{solucao}
\textbf{Passo 1 --- o crit\'erio de estabilidade.} Um matching \'e \emph{est\'avel} se n\~ao existe par bloqueante: um par candidato--curso que \emph{n\~ao} est\'a casado entre si, mas em que ambos prefeririam estar um com o outro do que com seu par atual em $\mu$. Note que o ju\'izo pedido \'e sobre a \emph{defini\c{c}\~ao} de estabilidade (exist\^encia de par bloqueante), e n\~ao sobre qual lado o algoritmo favorece.

\textbf{Passo 2 --- testar o par (Ana, Medicina).} Em $\mu$, Ana est\'a em Direito e Medicina est\'a com Bruno.
\begin{itemize}
\item \emph{Ana prefere Medicina?} Sim --- sua prefer\^encia \'e Medicina $\succeq$ Direito (estrita), e em $\mu$ ela est\'a em Direito.
\item \emph{Medicina prefere Ana?} Sim --- Medicina prioriza Ana sobre Bruno, e em $\mu$ ela tem Bruno.
\end{itemize}
Logo Ana e Medicina prefeririam, ambos, formar par entre si: $(\text{Ana},\text{Medicina})$ \'e par bloqueante.

\textbf{Passo 3 --- conclus\~ao sobre $\mu$.} Como existe par bloqueante, $\mu$ \textbf{n\~ao \'e est\'avel}. A afirma\c{c}\~ao est\'a correta. (Como contraste, o matching est\'avel seria Ana $\to$ Medicina e Bruno $\to$ Direito: a\'i Medicina tem sua candidata priorit\'aria e Direito tamb\'em, sem nenhum par bloqueante.)

\textbf{Passo 4 --- o que a afirma\c{c}\~ao \emph{n\~ao} alega.} A quest\~ao n\~ao fala em otimalidade-do-proponente nem em qual lado o algoritmo favorece; testa apenas a defini\c{c}\~ao de estabilidade via exist\^encia de um par bloqueante concreto. Por isso o julgamento \'e direto e independe de quem prop\~oe.

\textbf{Intui\c{c}\~ao econ\^omica:} a estabilidade \'e a condi\c{c}\~ao m\'inima para uma aloca\c{c}\~ao ``sobreviver'' a recursos e renegocia\c{c}\~oes. Se Ana, com prioridade suficiente, quer Medicina, e Medicina a prioriza sobre o aluno que ocupa a vaga, os dois t\^em incentivo a furar a aloca\c{c}\~ao proposta --- Ana recorreria e Medicina a aceitaria. Uma aloca\c{c}\~ao assim n\~ao se sustenta. \'E exatamente esse tipo de inconsist\^encia que a aceita\c{c}\~ao diferida elimina ao produzir um matching est\'avel.
\end{solucao}

\ifsolucao\else
\par\vspace{0.5em}
\begin{center}
\textbf{Quadro-resposta do Bloco 4 --- escreva a alternativa (A--E) ou V/F de cada item:}\\[0.6em]
\renewcommand{\arraystretch}{1.45}
\begin{tabular}{|>{\centering\arraybackslash}p{2.6cm}|>{\centering\arraybackslash}p{1.4cm}|>{\centering\arraybackslash}p{1.4cm}|>{\centering\arraybackslash}p{1.4cm}|}
  \hline
  \textbf{Item} & a) & b) & c) \\ \hline
  \textbf{Resposta} &  &  &  \\ \hline
\end{tabular}
\end{center}
{\small\itshape Atenção: apenas o quadro acima será considerado na correção dos itens objetivos deste bloco.}\par\vspace{0.6em}
\fi

\textbf{Questão 8.} (17 pontos) A certifica\c{c}\~ao CPA-20 da Anbima \'e exigida de profissionais que distribuem produtos de investimento em bancos e corretoras no Brasil. Suponha que ela funcione como \emph{sinal} \`a la Spence: o banco empregador n\~ao observa diretamente a produtividade do candidato, mas observa o n\'ivel de esfor\c{c}o de certifica\c{c}\~ao $e\ge 0$ (horas de estudo/qualifica\c{c}\~ao) registrado. H\'a dois tipos de candidato, com produtividades $\theta_L=1$ e $\theta_H=2$ (em unidades monet\'arias). Diferentemente do custo linear, aqui o custo de adquirir o sinal \'e \textbf{n\~ao-linear}, $c(e,\theta)=\dfrac{e^{2}}{2\theta}$ --- convexo em $e$ e mais barato, na margem, para o tipo alto, o que garante a condi\c{c}\~ao de \emph{single-crossing} (a taxa marginal de subs\-tit\-ui\c{c}\~ao $\partial c/\partial e=e/\theta$ decresce em $\theta$). Bancos competitivos pagam sal\'ario igual \`a produtividade esperada inferida do sinal (e das cren\c{c}as fora do equil\'ibrio). Metade dos candidatos \'e do tipo $H$.

\textbf{a)} (6 pontos) Considere primeiro o equil\'ibrio \emph{agrupador} (\emph{pooling}), em que ningu\'em sinaliza ($e=0$ para ambos os tipos) e o banco paga a todos a produtividade esperada. (i) Determine o sal\'ario agrupador. (ii) Esse equil\'ibrio \'e robusto? Verifique aplicando a l\'ogica do crit\'erio intuitivo de Cho--Kreps: usando o custo n\~ao-linear $c(e,\theta)=e^{2}/(2\theta)$, encontre o conjunto de sinais $e>0$ tais que um \emph{desvio} para esse $e$ seria estritamente lucrativo para o tipo $H$ mas n\~ao para o tipo $L$ (de modo que o banco, ao v\^e-lo, infira ``\'e $H$'' e pague $\theta_H$). Indique o desvio de \emph{menor custo} que quebra o agrupador.
\resolucaobox{5.5cm}
\begin{solucao}
\textbf{Passo 1 --- sal\'ario agrupador.} No \emph{pooling}, o banco n\~ao distingue tipos e paga a todos a produtividade esperada
\[
\bar\theta=\tfrac12\,\theta_L+\tfrac12\,\theta_H=\tfrac12\cdot 1+\tfrac12\cdot 2=1{,}5.
\]
Como ningu\'em sinaliza ($e=0$), o \emph{payoff} de cada tipo \'e $\bar\theta-c(0,\theta)=1{,}5$.

\textbf{Passo 2 --- a l\'ogica do desvio (Cho--Kreps).} Para testar a robustez, perguntamos: existe um sinal $e>0$ tal que, se o banco interpretasse ``quem fizer $e$ \'e $H$'' (pagando $\theta_H=2$), o tipo $H$ \emph{quereria} fazer esse desvio e o tipo $L$ \emph{n\~ao}? Se existe, a cren\c{c}a de que tal desvio s\'o pode vir de $H$ \'e a \'unica que sobrevive ao crit\'erio intuitivo, e o agrupador \emph{n\~ao} \'e robusto. Lembre que agora o custo \'e \emph{quadr\'atico}: $c(e,\theta_H)=e^{2}/4$ e $c(e,\theta_L)=e^{2}/2$.

\textbf{Passo 3 --- desvio lucrativo para $H$.} O tipo $H$ desvia para $e$ (recebendo $2$) em vez de ficar no pooling (recebendo $1{,}5$) se
\[
2-\frac{e^{2}}{2\theta_H}\;\ge\;1{,}5
\;\Longleftrightarrow\;
2-\frac{e^{2}}{4}\ge 1{,}5
\;\Longleftrightarrow\;
\frac{e^{2}}{4}\le 0{,}5
\;\Longleftrightarrow\;
e^{2}\le 2
\;\Longleftrightarrow\;
e\le \sqrt{2}.
\]

\textbf{Passo 4 --- desvio \emph{n\~ao} lucrativo para $L$.} O tipo $L$ \emph{n\~ao} quer mimetizar o desvio (mesmo recebendo $2$) se
\[
2-\frac{e^{2}}{2\theta_L}\;<\;1{,}5
\;\Longleftrightarrow\;
2-\frac{e^{2}}{2}<1{,}5
\;\Longleftrightarrow\;
\frac{e^{2}}{2}> 0{,}5
\;\Longleftrightarrow\;
e^{2}>1
\;\Longleftrightarrow\;
e>1.
\]
Combinando, qualquer
\[
e\in(1,\;\sqrt{2}\,]
\]
\'e um desvio que s\'o o tipo $H$ faz: lucrativo para ele, repelido pelo $L$. O \emph{single-crossing} (custo marginal $e/\theta$ menor para $H$) garante que essa janela \'e n\~ao-vazia, pois $1<\sqrt{2}\approx 1{,}41$. Logo o agrupador \textbf{n\~ao sobrevive} ao crit\'erio intuitivo. O desvio de \emph{menor custo} que o quebra \'e o piso da janela, $e=1$: nele, o tipo $L$ fica indiferente entre mimetizar ou n\~ao ($2-1/2=1{,}5$, igual ao pooling), e o tipo $H$ obt\'em $2-1/4=1{,}75>1{,}5$.

\textbf{Intui\c{c}\~ao econ\^omica:} o agrupador paga ao tipo alto menos do que ele vale ($1{,}5<2$) e ao tipo baixo mais ($1{,}5>1$). O tipo alto tem, ent\~ao, incentivo a ``gritar'' que \'e alto --- e o single-crossing lhe d\'a um grito que o tipo baixo n\~ao consegue (ou n\~ao quer) imitar. Com o custo convexo, o piso da janela sobe para $e=1$ (em vez do $0{,}5$ do caso linear), porque a primeira hora de estudo agora ``machuca'' menos na margem; ainda assim, basta um sinal modesto para romper o disfarce coletivo. Por isso o pooling, embora atraente \`a primeira vista, n\~ao se sustenta sob cren\c{c}as razo\'aveis fora do equil\'ibrio.
\rubricalabel
\begin{itemize}
\item Calcula o sal\'ario agrupador $\bar\theta=1{,}5$ pela produtividade esperada: 1{,}5 pt.
\item Reconhece que testar robustez exige um desvio lucrativo para $H$ e n\~ao para $L$ (l\'ogica Cho--Kreps): 1 pt.
\item Deriva, com o custo \emph{quadr\'atico}, a condi\c{c}\~ao de desvio lucrativo para $H$ ($e\le\sqrt{2}$): 1{,}25 pt.
\item Deriva a condi\c{c}\~ao de n\~ao-mimetiza\c{c}\~ao do $L$ ($e>1$): 1{,}25 pt.
\item Conclui a janela de quebra $e\in(1,\,\sqrt{2}\,]$ e identifica o desvio de menor custo $e=1$: 1 pt.
\item Usar o custo linear $e/\theta$ em vez do quadr\'atico (obtendo a janela $(0{,}5,\,1]$): teto de 3 pts no item.
\item Inverter os pap\'eis (pedir desvio lucrativo para $L$ e n\~ao para $H$): teto de 3 pts no item.
\item Erro alg\'ebrico isolado (ex.: troca de sinal numa desigualdade ou esquecer a raiz) com a montagem correta: penalidade leve de $-0{,}5$ pt, n\~ao zera o item.
\end{itemize}
\end{solucao}

\textbf{b)} (6 pontos) Como o agrupador n\~ao \'e robusto, a previs\~ao recai sobre os equil\'ibrios \emph{separadores}, em que o tipo $L$ escolhe $e=0$ (recebe $\theta_L$) e o tipo $H$ escolhe $e_H>0$ (recebe $\theta_H$). Mantenha o custo n\~ao-linear $c(e,\theta)=e^{2}/(2\theta)$. (i) Como \emph{subproduto} das duas restri\c{c}\~oes de compatibilidade de incentivo, derive a faixa $[\underline e,\,\overline e]$ de sinais $e_H$ que sustentam a separa\c{c}\~ao e calcule-a numericamente. (ii) Aplicando o crit\'erio de Riley (menor custo), identifique o sinal $e^{*}$, o custo de certifica\c{c}\~ao incorrido pelo tipo $H$ e verifique explicitamente que ambas as restri\c{c}\~oes s\~ao satisfeitas.
\resolucaobox{5.5cm}
\begin{solucao}
\textbf{Passo 1 --- as duas restri\c{c}\~oes de incentivo.} O \emph{payoff} de um tipo $\theta$ com sinal $e$ e sal\'ario $w$ \'e $w-c(e,\theta)=w-e^{2}/(2\theta)$. Para separar:
\[
\underbrace{\theta_L-0\;\ge\;\theta_H-\frac{e_H^{2}}{2\theta_L}}_{L\ \text{n\~ao imita }H}
\;\Longleftrightarrow\;
\frac{e_H^{2}}{2\theta_L}\ge\theta_H-\theta_L
\;\Longleftrightarrow\;
e_H\ge\sqrt{2\theta_L(\theta_H-\theta_L)}\equiv\underline e,
\]
\[
\underbrace{\theta_H-\frac{e_H^{2}}{2\theta_H}\;\ge\;\theta_L-0}_{H\ \text{prefere sinalizar}}
\;\Longleftrightarrow\;
\frac{e_H^{2}}{2\theta_H}\le\theta_H-\theta_L
\;\Longleftrightarrow\;
e_H\le\sqrt{2\theta_H(\theta_H-\theta_L)}\equiv\overline e.
\]
A faixa separadora \'e $e_H\in[\underline e,\,\overline e]=\big[\sqrt{2\theta_L(\theta_H-\theta_L)},\,\sqrt{2\theta_H(\theta_H-\theta_L)}\,\big]$, n\~ao-vazia pelo single-crossing ($\theta_H>\theta_L>0\Rightarrow\underline e<\overline e$).

\textbf{Passo 2 --- valores num\'ericos.} Com $\theta_L=1$, $\theta_H=2$ (logo $\theta_H-\theta_L=1$):
\[
\underline e=\sqrt{2\cdot 1\cdot 1}=\sqrt{2}\approx 1{,}41,\qquad \overline e=\sqrt{2\cdot 2\cdot 1}=\sqrt{4}=2,\qquad e_H\in[\sqrt{2},\,2].
\]

\textbf{Passo 3 --- equil\'ibrio de Riley (menor custo).} Entre todos os separadores, o de menor custo escolhe o \emph{menor} sinal que ainda dissuade o tipo $L$, isto \'e, $e^{*}=\underline e=\sqrt{2}$. Qualquer $e_H<\sqrt{2}$ violaria a IC do tipo $L$. O custo incorrido pelo tipo $H$ \'e
\[
c(e^{*},\theta_H)=\frac{(e^{*})^{2}}{2\theta_H}=\frac{(\sqrt{2})^{2}}{2\cdot 2}=\frac{2}{4}=0{,}5,
\]
e seu \emph{payoff} l\'iquido \'e $2-0{,}5=1{,}5$.

\textbf{Passo 4 --- verifica\c{c}\~ao das duas ICs em $e^{*}=\sqrt{2}$.}
\begin{itemize}
\item Tipo $L$ n\~ao imita: ficar como $L$ rende $\theta_L-0=1$; imitar $H$ renderia $\theta_H-\dfrac{(e^{*})^{2}}{2\theta_L}=2-\dfrac{2}{2}=2-1=1$. Como $1\ge 1$, a IC do tipo $L$ vale (\emph{binding}, como esperado em Riley).
\item Tipo $H$ sinaliza: sinalizar rende $\theta_H-\dfrac{(e^{*})^{2}}{2\theta_H}=2-0{,}5=1{,}5$; desistir do sinal e ser tratado como $L$ renderia $\theta_L=1$. Como $1{,}5\ge 1$, a IC do tipo $H$ vale com folga.
\end{itemize}
Ambas satisfeitas: $e^{*}=\sqrt{2}$ \'e o separador v\'alido de menor custo.

\textbf{Intui\c{c}\~ao econ\^omica:} note que o \emph{payoff} l\'iquido do tipo $H$ em Riley ($1{,}5$) coincide com o que ele teria no agrupador --- a sinaliza\c{c}\~ao apenas consome a renda que o pooling lhe daria de gra\c{c}a. Com o custo convexo, o sinal $e^{*}=\sqrt{2}$ \'e maior do que o do caso linear, mas o \emph{custo monet\'ario} queimado ($0{,}5$) sai o mesmo, porque o que ancora a separa\c{c}\~ao \'e o \emph{custo} que o tipo $L$ pagaria ao imitar, n\~ao o n\'ivel bruto do sinal. O tipo $L$ \'e quem ``ancora'' o sinal m\'inimo: a certifica\c{c}\~ao tem de doer o suficiente para que o profissional menos produtivo n\~ao ache que vale a pena estudar para se passar por mais qualificado. No menor sinal que cumpre isso, o tipo $L$ fica indiferente entre imitar ou n\~ao (IC binding), e o tipo $H$ ainda se distingue. Riley seleciona exatamente esse ponto: a menor quantidade de recurso queimada compat\'ivel com a separa\c{c}\~ao.
\rubricalabel
\begin{itemize}
\item Monta corretamente a IC do tipo $L$ \emph{com custo quadr\'atico} e deriva $\underline e=\sqrt{2\theta_L(\theta_H-\theta_L)}$: 1{,}25 pt.
\item Monta corretamente a IC do tipo $H$ \emph{com custo quadr\'atico} e deriva $\overline e=\sqrt{2\theta_H(\theta_H-\theta_L)}$: 1{,}25 pt.
\item Calcula numericamente $\underline e=\sqrt{2}\approx 1{,}41$ e $\overline e=2$: 1 pt (0{,}5 cada).
\item Identifica o equil\'ibrio de Riley $e^{*}=\underline e=\sqrt{2}$ e o custo $c(e^{*},\theta_H)=0{,}5$: 1{,}25 pt.
\item Verifica a IC do tipo $L$ ($1\ge 1$, binding) e a do tipo $H$ ($1{,}5\ge 1$): 1{,}25 pt.
\item Escolher $e^{*}=\overline e=2$ chamando-o de ``menor custo'': perde os 1{,}25 pt da identifica\c{c}\~ao de Riley, mas mant\'em os pontos das verifica\c{c}\~oes se coerentes.
\item Usar o custo linear (obtendo faixa $[1,\,2]$ sem raiz): perde os 2{,}5 pts das duas deriva\c{c}\~oes das ICs; mant\'em parcialmente os demais se coerentes.
\item Erro aritm\'etico isolado com m\'etodo correto: penalidade leve de $-0{,}5$ pt.
\end{itemize}
\end{solucao}

\textbf{c)} (5 pontos) O sinal n\~ao aumenta a produtividade de ningu\'em: \'e puro custo social. Use o custo $c(e,\theta)=e^{2}/(2\theta)$. (i) Calcule o desperd\'icio social esperado por candidato no equil\'ibrio separador de \emph{menor custo} ($e^{*}=\sqrt{2}$) e no de \emph{maior custo} ($e_H=\overline e=2$), lembrando que metade dos candidatos \'e do tipo $H$. (ii) Indique o desperd\'icio social da aloca\c{c}\~ao \emph{agrupadora}. (iii) Ordene as tr\^es situa\c{c}\~oes em efici\^encia (desperd\'icio de recursos) e explique por que, apesar de o agrupador ser o mais eficiente, ele n\~ao \'e a previs\~ao do modelo.
\resolucaobox{5.5cm}
\begin{solucao}
\textbf{Passo 1 --- desperd\'icio no separador de menor custo.} O \'unico recurso queimado \'e o custo de certifica\c{c}\~ao do tipo $H$, $c(e^{*},\theta_H)=\dfrac{(\sqrt{2})^{2}}{2\cdot 2}=\dfrac{2}{4}=0{,}5$, incorrido pela metade $H$ dos candidatos. O desperd\'icio social esperado por candidato \'e
\[
D_{\min}=\tfrac12\cdot c(\sqrt{2},\theta_H)+\tfrac12\cdot 0=\tfrac12\cdot 0{,}5=0{,}25.
\]

\textbf{Passo 2 --- desperd\'icio no separador de maior custo.} Com $e_H=2$, o custo do tipo $H$ \'e $c(2,\theta_H)=\dfrac{2^{2}}{2\cdot 2}=\dfrac{4}{4}=1$. Ent\~ao
\[
D_{\max}=\tfrac12\cdot 1+\tfrac12\cdot 0=0{,}5.
\]
Na faixa separadora, o desperd\'icio por candidato varia em $[0{,}25,\,0{,}5]$: quanto mais alto o sinal exigido, mais recursos se queimam para transmitir a \emph{mesma} informa\c{c}\~ao.

\textbf{Passo 3 --- aloca\c{c}\~ao agrupadora.} No \emph{pooling}, ningu\'em sinaliza ($e=0$ para todos), de modo que nenhum recurso \'e queimado: $D_{\text{pool}}=0$.

\textbf{Passo 4 --- ordena\c{c}\~ao e previs\~ao do modelo.} Em efici\^encia (do menos ao mais desperdi\c{c}ador):
\[
D_{\text{pool}}=0\;<\;D_{\min}=0{,}25\;<\;D_{\max}=0{,}5.
\]
O agrupador \'e o mais eficiente --- a sinaliza\c{c}\~ao nada acrescenta \`a produtividade, ent\~ao queim\'a-la \'e perda l\'iquida. Mas, como mostrado no sub-item (a), o agrupador \textbf{n\~ao \'e robusto}: o tipo $H$ disp\~oe de um desvio (qualquer $e\in(1,\,\sqrt{2}\,]$) que sinaliza credivelmente seu tipo --- lucrativo para ele, repelido pelo tipo $L$ via single-crossing --- e que sobrevive ao crit\'erio intuitivo de Cho--Kreps. Por isso o pooling desmorona, e a previs\~ao robusta \'e o equil\'ibrio separador, com desperd\'icio estritamente positivo. Entre os separadores, Riley seleciona o de menor custo, $D_{\min}=0{,}25$.

\textbf{Intui\c{c}\~ao econ\^omica:} este \'e o paradoxo central de Spence --- a sinaliza\c{c}\~ao \'e socialmente \emph{custosa} sem ser socialmente \emph{produtiva}. A economia gasta recursos reais (horas de estudo, taxas de certifica\c{c}\~ao) apenas para reembaralhar informa\c{c}\~ao que j\'a existe, sem produzir nada a mais. O agrupador eliminaria esse desperd\'icio, mas \'e inst\'avel porque o tipo alto sempre tem como, e interesse em, gritar que \'e alto. O separador de menor custo \'e o ``menos pior'': sustenta-se e queima o m\'inimo de recursos compat\'ivel com a separa\c{c}\~ao.
\rubricalabel
\begin{itemize}
\item Calcula $D_{\min}=0{,}25$ com $c(\sqrt{2},\theta_H)=0{,}5$, ponderando pela fra\c{c}\~ao $\tfrac12$ de tipos $H$: 1{,}5 pt.
\item Calcula $D_{\max}=0{,}5$ com $c(2,\theta_H)=1$: 1 pt.
\item Identifica $D_{\text{pool}}=0$ (ningu\'em sinaliza): 1 pt.
\item Ordena corretamente $D_{\text{pool}}<D_{\min}<D_{\max}$: 0{,}5 pt.
\item Explica por que o agrupador n\~ao \'e a previs\~ao (n\~ao-robusto: tipo $H$ desvia e sinaliza credivelmente, single-crossing / Cho--Kreps): 1 pt.
\item Esquecer a pondera\c{c}\~ao $\tfrac12$ (reportar $0{,}5$ e $1$ como desperd\'icios ``por candidato''): perde 0{,}5 pt, desde que a ordena\c{c}\~ao permane\c{c}a coerente.
\item Erro aritm\'etico isolado com m\'etodo correto: penalidade leve de $-0{,}5$ pt.
\end{itemize}
\end{solucao}
